\documentclass[a4paper, 12pt, oneside, openany]{article}
\frenchspacing
\usepackage[left=3cm,right=1.5cm,top=2cm,bottom=2cm]{geometry}
\pagestyle{plain}

% ✅ Кодировка и шрифты (ПРАВИЛЬНЫЙ ПОРЯДОК)
\usepackage[T2A]{fontenc}      % ✅ Только T2A для русского!
\usepackage[utf8]{inputenc}
\usepackage{lmodern}           % ✅ Масштабируемые шрифты
\usepackage[russian]{babel}

% ✅ Остальные пакеты
\usepackage{float}
\usepackage{flafter}
\usepackage{wrapfig}
\usepackage[pdftex,unicode]{hyperref}
% ... гиперссылки настройки ...

% ✅ Microtype (после шрифтов, с final)
\usepackage[expansion=false, protrusion=false]{microtype}  % ✅ Отключить всё проблемное

\usepackage{amsmath, amssymb}
\usepackage[warn]{mathtext}

\usepackage{pgfplots}
\pgfplotsset{compat=newest}
\usepackage{tikz}              % ✅ ДОБАВИТЬ СЮДА
\usetikzlibrary{shapes,arrows,positioning,calc}


%% Реализация библиографии пакетами biblatex и biblatex-gost с использованием движка biber - не работает в overleaf из-за долгой компиляции в бесплатном режиме. В других приложениях можно брать %%%
%
%\usepackage{csquotes} % Пакет для оформления сложных блоков цитирования (biblatex рекомендует его подключать)
%\usepackage[%
%backend=biber,                                      % Движок
%bibencoding=utf8,                                   % Кодировка bib-файла
%sorting=none,                                       % Настройка сортировки списка литературы
%style=gost-numeric,                                 % Стиль цитирования и библиографии по ГОСТ
%language=auto,                                      % Язык для каждой библиографической записи задается отдельно
%autolang=other,                                     % Поддержка многоязычной библиографии
%sortcites=true,                                     % Если в квадратных скобках несколько ссылок, то отображаться будут отсортированно
%movenames=false,                                    % Не перемещать имена, они всегда в начале библиографической записи
% maxnames=5,                                         % Максимальное отображаемое число авторов
% minnames=3,                                         % До скольки сокращать число авторов, если их больше максимума
% doi=false,                                          % Не отображать ссылки на DOI
% isbn=false,                                         % Не показывать ISBN, ISSN, ISRN
% ]{biblatex}[2016/09/17]
% \DeclareDelimFormat{bibinitdelim}{}                 % Убираем пробел между инициалами (Иванов И.И. вместо Иванов И. И.)
% \addbibresource{biba.bib}                           % Определяем файл с библиографией
%
% %%% Скрипт, который автоматически подбирает язык (и, следовательно, формат) для каждой библиографической записи %%%
% %%% Если в названии работы есть кириллица - меняем значение поля langid на russian %%%
%%%% Все оставшиеся пустые места в поле langid заменяем на english %%%
%
%\DeclareSourcemap{
%   \maps[datatype=bibtex]{
%     \map{
%         \step[fieldsource=title, match=\regexp{^\P{Cyrillic}*\p{Cyrillic}.*}, final]
%         \step[fieldset=langid, fieldvalue={russian}]
%     }
%     \map{
%         \step[fieldset=langid, fieldvalue={english}]
%     }
%   }
% }


%%% Задаем параметры оформления рисунков и таблиц %%%

\usepackage{graphicx, caption, subcaption}  % Подгружаем пакеты для работы с графикой и настройки подписей
\graphicspath{{images/}} % Определяем папку с рисунками
\RequirePackage{caption}
\DeclareCaptionLabelSeparator{defffis}{ - }
\captionsetup[figure]{font=small,  name=Рисунок, justification=centering, labelsep=defffis, labelfont=bf} % Задаем параметры подписей к рисункам: маленький шрифт (в данном случае 12pt), ширина равна ширине текста, полнотекстовая надпись "Рисунок", выравнивание по центру
\captionsetup[subfigure]{font=small} % Индексы подрисунков а), б) и так далее тоже шрифтом 12pt (по умолчанию делает еще меньше)
\captionsetup[table]{singlelinecheck=false,font=small,width=\textwidth,justification=justified, labelsep=defffis, labelfont=bf} % Задаем параметры подписей к таблицам: запрещаем переносы, маленький шрифт (в данном случае 12pt), ширина равна ширине текста, выравнивание по ширине
% \captiondelim{-} % Разделителем между номером рисунка/таблицы и текстом в подписи является длинное тире. Не работает
\setkeys{Gin}{width=\textwidth} % По умолчанию размер всех добавляемых рисунков будет подгоняться под ширину текста
\renewcommand{\thesubfigure}{\asbuk{subfigure}} % Нумерация подрисунков строчными буквами кириллицы
%\setlength{\abovecaptionskip}{0pt} % Отбивка над подписью - тут не меняем
%\setlength{\belowcaptionskip}{0pt} % Отбивка под подписью - тут не меняем
\usepackage[section]{placeins} % Объекты типа float (рисунки/таблицы) не вылезают за границы секциии, в которой они объявлены

\usepackage{longtable} %можно делать многостраничные таблицы
\usepackage{rotating} %можно поворачивать объекты на заданный угол \rotatebox{90}{}
\usepackage{multirow}
\usepackage{color,colortbl}
\definecolor{darkishgreen}{RGB}{39,203,22}
\definecolor{LightCyan}{rgb}{0.88,1,1}
\definecolor{Gray}{gray}{0.9}
\definecolor{lightRed}{RGB}{230,170,150}
\definecolor{modRed}{RGB}{230,82,90}
\definecolor{strongRed}{RGB}{230,6,6}
\setlength{\parindent}{1.25cm} % Устанавливает отступ первой строки 1.25 см
\usepackage{enumitem}
\renewcommand*{\labelitemi}{\normalfont{-}}        % В ненумерованных списках для пунктов используем короткое тире
\makeatletter
    \AddEnumerateCounter{\asbuk}{\russian@alph}     % Объясняем пакету enumitem, как использовать asbuk
\makeatother
\renewcommand{\labelenumii}{\asbuk{enumii})}        % Кириллица для второго уровня нумерации
\renewcommand{\labelenumiii}{\arabic{enumiii})}     % Арабские цифры для третьего уровня нумерации
\setlist{noitemsep, leftmargin=*}                   % Убираем интервалы между пунками одного уровня в списке
% \setlist[1]{labelindent=\parindent}                 % Отступ у пунктов списка равен абзацному отступу
% \setlist[2]{leftmargin=\parindent}                  % Плюс еще один такой же отступ для следующего уровня
% \setlist[3]{leftmargin=\parindent}                  % И еще один для третьего уровня

\setlist[itemize]{leftmargin=-0.15cm, itemindent=1.25cm} % Для всех itemize
\setlist[enumerate]{leftmargin=0.25cm, itemindent=1.25cm} % Для всех enumerate
\usepackage{pdflscape}
\linespread{1.25}
\clubpenalty=10000                                  % Запрещаем разрыв страницы после первой строки абзаца
\widowpenalty=10000                                 % Запрещаем разрыв страницы после последней строки абзаца
\brokenpenalty=4991                                 % Ограничение на разрыв страницы, если строка заканчивается переносом

\usepackage{appendix} % Для удобного создания приложений
\usepackage{tocloft}
\renewcommand{\cftsecleader}{\cftdotfill{\cftdotsep}}
% \newcommand{\intro}[1]{
%     \stepcounter{section}
%     \section*{\centeringПриложение \Asbuk{section}}
%     \begin{center}
%         \bf{#1}
%     \end{center}
%     \markboth{\MakeUppercase{#1}}{}
%     \addcontentsline{toc}{section}{Приложение \Asbuk{section}. #1}
% }
\usepackage{mdframed}
\usepackage{xcolor} % Для цвета рамки (опционально)
\usepackage{listings}
\DeclareCaptionFont{white}{\color{white}} %% это сделает текст заголовка белым
%% код ниже нарисует серую рамочку вокруг заголовка кода.
\DeclareCaptionFormat{listing}{\colorbox{gray}{\parbox{\textwidth}{#1#2#3}}}
\captionsetup[lstlisting]{format=listing,labelfont=white,textfont=white, singlelinecheck=false,font=small,width=\textwidth,justification=justified, labelsep=defffis}

\lstset{ 
language=Python,
basicstyle=\ttfamily\small, % размер и начертание шрифта для подсветки кода
numbers=left,               % где поставить нумерацию строк (слева\справа)
numberstyle=\footnotesize,           % размер шрифта для номеров строк
stepnumber=1,                   % размер шага между двумя номерами строк
numbersep=5pt,                % как далеко отстоят номера строк от подсвечиваемого кода
showspaces=false,            % показывать или нет пробелы специальными отступами
showstringspaces=false,      % показывать или нет пробелы в строках
showtabs=false,             % показывать или нет табуляцию в строках
tabsize=2,                 % размер табуляции по умолчанию равен 2 пробелам
captionpos=t,              % позиция заголовка вверху [t] или внизу [b] 
breaklines=true,           % автоматически переносить строки (да\нет)
breakatwhitespace=false, % переносить строки только если есть пробел
% escapeinside={\%}{*)},   % если нужно добавить комментарии в коде
escapeinside={(*@}{@*)},  % ✅ Разные открывающий и закрывающий
columns=fullflexible,
flexiblecolumns,
extendedchars=true,
inputencoding=utf8,
  % Для кириллицы в комментариях:
  literate={а}{{\cyra}}1 {б}{{\cyrb}}1 {в}{{\cyrv}}1 {г}{{\cyrg}}1 
           {д}{{\cyrd}}1 {е}{{\cyre}}1 {ё}{{\cyryo}}1 {ж}{{\cyrzh}}1
           {з}{{\cyrz}}1 {и}{{\cyri}}1 {й}{{\cyrishrt}}1 {к}{{\cyrk}}1
           {л}{{\cyrl}}1 {м}{{\cyrm}}1 {н}{{\cyrn}}1 {о}{{\cyro}}1
           {п}{{\cyrp}}1 {р}{{\cyrr}}1 {с}{{\cyrs}}1 {т}{{\cyrt}}1
           {у}{{\cyru}}1 {ф}{{\cyrf}}1 {х}{{\cyrh}}1 {ц}{{\cyrc}}1
           {ч}{{\cyrch}}1 {ш}{{\cyrsh}}1 {щ}{{\cyrshch}}1 {ъ}{{\cyrhrds}}1
           {ы}{{\cyrery}}1 {ь}{{\cyrsftsn}}1 {э}{{\cyre}}1 {ю}{{\cyryu}}1
           {я}{{\cyrya}}1
           {А}{{\CyrA}}1 {Б}{{\CyrB}}1 {В}{{\CyrV}}1 {Г}{{\CyrG}}1
           {Д}{{\CyrD}}1 {Е}{{\CyrE}}1 {Ё}{{\CyrYo}}1 {Ж}{{\CyrZh}}1
           {З}{{\CyrZ}}1 {И}{{\CyrI}}1 {Й}{{\CyrIi}}1 {К}{{\CyrK}}1
           {Л}{{\CyrL}}1 {М}{{\CyrM}}1 {Н}{{\CyrN}}1 {О}{{\CyrO}}1
           {П}{{\CyrP}}1 {Р}{{\CyrR}}1 {С}{{\CyrS}}1 {Т}{{\CyrT}}1
           {У}{{\CyrU}}1 {Ф}{{\CyrF}}1 {Х}{{\CyrH}}1 {Ц}{{\CyrTs}}1
           {Ч}{{\CyrCh}}1 {Ш}{{\CyrSh}}1 {Щ}{{\CyrShch}}1 {Ъ}{{\CyrHrds}}1
           {Ы}{{\CyrEry}}1 {Ь}{{\CyrSftsn}}1 {Э}{{\CyrE}}1 {Ю}{{\CyrYu}}1
           {Я}{{\CyrYa}}1,
keywordstyle=\color{blue}\bfseries,
identifierstyle=\color{black},
morecomment=[l]{\#},  % ✅ Экранировать символ #
commentstyle=\color{green}\ttfamily,
stringstyle=\color{red}\ttfamily,
morestring=[b]",
keepspaces=true,
% escapechar=\%,
texcl=false,
backgroundcolor=\color{white}, % цвет фона подсветки - используем \usepackage{color}
frame=single              % рисовать рамку вокруг кода
}
% \usepackage{minted}
% \usemintedstyle{colorful}
% \DeclareCaptionFont{white}{\color{white}} %% это сделает текст заголовка белым
% %% код ниже нарисует серую рамочку вокруг заголовка кода.
% \DeclareCaptionFormat{minted}{\colorbox{gray}{\parbox{\textwidth}{#1#2#3}}}
% \captionsetup[minted]{format=minted,labelfont=white,textfont=white, singlelinecheck=false,font=small,width=\textwidth,justification=justified, labelsep=defffis}


\begin{document}
\begin{titlepage} % начало титульной страницы
\begin{center} % включить выравнивание по центру
Федеральное государственно автономное образовательное учреждение\\ высшего образования\\
<<Московский физико-технический институт (государственный университет)>>\\
Центр дополнительного, дополнительного профессионального\\и онлайн-образования <<Пуск>>\\ [2cm]
\end{center} % закончить выравнивание по центру
\par
\textbf{Направление подготовки:} 01.04.02 Прикладная математика и информатика

\textbf{Направленность (профиль) подготовки:} Науки о данных

\vspace{2cm}

\begin{center} % включить выравнивание по центру
{\huge\textbf{Наименование темы ВКР}}\\
(магистерская диссертация)\\[2cm]
\end{center} % закончить выравнивание по центру
\flushright {
\begin{minipage}{0.5\textwidth} % начало маленькой врезки в половину ширины текста
\begin{flushleft} % выровнять её содержимое по левому краю
\textbf{Студент:}\\
ФИО\\
\hrulefill
\parskip=-5mm
\begin{center}\footnotesize\textit{(подпись студента)}\end{center}
\parskip=3mm

\textbf{Научный руководитель:}
\par
ФИО\\
\hrulefill
\parskip=-5mm
\begin{center} \footnotesize\emph{(подпись научного руководителя)}\end{center}
\parskip=3mm
\textbf{Консультант:}\\
ФИО\\
\hrulefill 
\parskip=-5mm
\begin{center} \footnotesize\emph{(подпись консультанта)}\end{center}

\end{flushleft} % конец выравнивания по левому краю
\end{minipage} % конец врезки
}
\vfill % заполнить всё доступное ниже пространство
\begin{center}{Москва \the\year}\end{center} % вывести дату
\thispagestyle{empty} % не нумеровать страницу
\end{titlepage} % конец титульной страницы

\newpage
\setcounter{page}{2}
\section{Аннотация}
Магистерская диссертация посвящена разработке и экспериментальной валидации кросс-модальной функции потерь, применяемой в режиме файнтюнинга к двум опорным архитектурам генерации говорящих лиц - Wav2Lip и SadTalker - для повышения эмоциональной согласованности между речевым сигналом и мимикой на подмножестве из четырёх эмоций (\textit{happy}, \textit{sad}, \textit{angry}, \textit{disgust}) датасета RAVDESS. Опорные архитектуры ориентированы на синхронизацию артикуляции и не контролируют соответствие эмоционального содержания речи и видеоряда. Существующие подходы к управлению эмоциональной выразительностью (EmoTalk, EmoTalker, EmotiveTalk) предполагают замену архитектуры или использование явных меток на этапе генерации, что ограничивает их применимость к уже развёрнутым моделям.

В работе предложена кросс-модальная функция потерь, реализуемая как надстройка над функцией реконструкции L1 и связывающая выходы предобученных и зафиксированных аудио- и видеоэнкодеров эмоций. Эмпирически установленная победная конфигурация совмещает перекрёстную энтропию на категориальных метках эмоций и KL-дивергенцию между логитами видеоэнкодера и сглаженным распределением аудиоэнкодера ($T=2{,}0$). Подход не требует явной разметки эмоций на этапе инференса и не предполагает изменения архитектуры базовой модели.

Экспериментальная валидация проведена на датасете RAVDESS с разбиением по актёрам без пересечений (544/96/96 сэмплов в выборках обучения, валидации и тестирования). Аудиоэнкодер реализован файнтюнингом модели \texttt{superb/wav2vec2-base-superb-er} и достигает на тестовой выборке макро-F1 0,739; видеоэнкодер реализован файнтюнингом модели \texttt{facebook/timesformer-base-finetuned-k400} (восемь кадров на сэмпл) и достигает макро-F1 0,876. Проведена пятиконфигурационная абляция функции потерь (рекон\-струкционный baseline; косинусное сближение; перекрёстная энтропия; перекрёстная энтропия совместно с косинусом; перекрёстная энтропия совместно с KL-дистилляцией). Победившая по F1 на валидации конфигурация (CE+KL) использована для основного файнтюнинга Wav2Lip и SadTalker с поиском весового коэффициента. Для снижения риска цикличности оценки - обусловленной использованием внутренних энкодеров одновременно в функции потерь и в первичной метрике - проведена независимая внешняя оценка с применением двух классификаторов (\texttt{motheecreator/vit-Facial-Expression-Recognition}, \texttt{Rajaram1996/FacialEmoRecog}), не задействованных ни на одном этапе обучения и отбора моделей; отбор внешних моделей выполнен по итогам скрининга пяти кандидатов на чистом RAVDESS.

На тестовой выборке файнтюнинг Wav2Lip с предложенной функцией потерь повышает F1-меру эмоциональной согласованности с 0,611 до 0,718 ($\Delta\text{F1}=+10{,}7$~\%) при снижении LSE-C с 0,242 до 0,227 ($\Delta\text{LSE-C}=-6{,}1$~\%); направление эффекта согласовано с обоими внешними классификаторами. Магнитуда эффекта по F1 удовлетворяет порогу гипотезы H1, однако статистическая значимость по критерию Макнемара на $n=96$ не достигнута ($p = 0{,}170$); снижение LSE-C превышает заявленный допуск в 2~\%, но не является статистически значимым ($p = 0{,}611$). На рендеренной выборке SadTalker ($n=24$ на конфигурацию) улучшения по внешним классификаторам не зафиксировано, что не позволяет считать гипотезу H2 подтверждённой в текущем эксперименте. Заявленный вклад ограничен двумя опорными архитектурами, четырьмя эмоциями и студийным датасетом RAVDESS; обобщение на иные архитектуры, расширенное эмоциональное пространство и in-the-wild видео выходит за пределы экспериментально подтверждённого диапазона.

\textbf{Ключевые слова:} генерация говорящих лиц, кросс-модальное обучение, эмоциональная согласованность, файнтюнинг, KL-дистилляция, внешняя валидация, Wav2Lip, SadTalker, RAVDESS.

\newpage
% \renewcommand{\contentsname}{Оглавление}
\tableofcontents  % содержание
\suppressfloats
\newpage

\section{Введение}

\subsection{Актуальность исследования}

Технологии генерации говорящих лиц применяются в цифровых ассистентах, образовательных платформах, телемедицине и игровой индустрии. По прогнозу Grand View Research, мировой рынок цифровых аватаров возрастёт с 18 миллиардов долларов США в 2023 году до 270 миллиардов долларов США к 2030 году [Grand View Research, 2024], что обусловливает повышение требований к естественности их эмоционального поведения.

Современные архитектуры аудио-управляемого синтеза говорящих лиц достигли высоких показателей точности артикуляции: Wav2Lip [Prajwal et al., 2020] обеспечивает точность синхронизации губ свыше 97~\%, SadTalker [Zhang et al., 2023] реализует генерацию трёхмерной геометрии лица и движений головы. Однако перечисленные модели ориентированы на синхронизацию артикуляции и не предусматривают явного контроля соответствия эмоционального содержания речи и мимики, вследствие чего при подаче эмоционально насыщенного сигнала генерируемое видео сохраняет нейтральное выражение. Согласно перцепционным исследованиям, рассогласование голосового и мимического каналов снижает точность распознавания эмоций пользователями и оценку реализма аватара [Suma et al., 2023; Zhou et al., 2025].

В период с 2022 по 2025 год опубликован ряд работ, в которых эмоционально-обусловленное управление является существенным компонентом метода: EmoTalk [Peng et al., 2023], EmoTalker [Shen et al., 2024], EmotiveTalk [Wang et al., 2025], EAT [Wang et al., 2023]. Перечисленные подходы либо требуют замены опорной архитектуры (диффузионный или специализированный backbone), либо предполагают явную разметку эмоций на этапе инференса. Возможность повышения эмоциональной согласованности при сохранении уже развёрнутых опорных архитектур и без явных меток на этапе генерации в литературе системно не исследована. Сужение теоретического пробела до указанной постановки определяет направление настоящего исследования.

\subsection{Объект и предмет исследования}

\textbf{Объект исследования} - модели аудио-управляемой генерации говорящих лиц Wav2Lip и SadTalker.

\textbf{Предмет исследования} - кросс-модальная функция потерь, связывающая эмоциональные представления речевого сигнала и видеоряда в едином векторном пространстве.

\subsection{Цель и задачи исследования}

\textbf{Цель исследования} - разработка и экспериментальная валидация кросс-модальной функции потерь, обеспечивающей статистически значимое улучшение F1-меры эмоциональной согласованности между речевым сигналом и мимикой при файнтюнинге опорных архитектур Wav2Lip и SadTalker на подмножестве из четырёх эмоций (\textit{happy}, \textit{sad}, \textit{angry}, \textit{disgust}) датасета RAVDESS без явной разметки эмоций на этапе генерации и без модификации архитектуры опорных моделей.

Для достижения поставленной цели решаются следующие задачи:
\begin{enumerate}
\item теоретический анализ существующих подходов к эмоционально-обусловленному синтезу говорящих лиц и формулировка теоретического пробела;
\item файнтюнинг аудио- и видеоэнкодеров эмоций на датасете RAVDESS как опорных моделей для извлечения эмоциональных представлений из соответствующих модальностей;
\item формализация кросс-модальной функции потерь и её абляционный анализ в пятиконфигурационном эксперименте на Wav2Lip с целью эмпирического выбора компонент (косинусное сближение, перекрёстная энтропия, KL-дистилляция и их сочетания);
\item файнтюнинг моделей Wav2Lip и SadTalker с применением победившей в абляции функции потерь и поиск весового коэффициента, обеспечивающего максимум F1 в пределах допустимой деградации реконструкции;
\item сравнительный анализ с базовыми моделями на тестовой выборке с проверкой статистической значимости (критерий Макнемара, парный $t$-критерий) и независимой внешней оценкой эмоциональной согласованности с использованием классификаторов, не задействованных в обучающей сигнальной цепи.
\end{enumerate}

\subsection{Гипотезы исследования}

\textbf{Гипотеза H1 (основная).} Включение предложенной кросс-модальной функции потерь в процедуру файнтюнинга Wav2Lip на четырёх эмоциях RAVDESS приводит к улучшению F1-меры эмоциональной согласованности не менее чем на 10~\% по сравнению с базовой моделью при условии, что абсолютное изменение метрики LSE-C не превышает 2~\% от значения базовой модели. \textit{Критерии подтверждения:} $\Delta\text{F1} \geq +10$~\% ($p < 0{,}05$, критерий Макнемара); $|\Delta\text{LSE-C}| \leq 2$~\% либо изменение LSE-C статистически незначимо по парному $t$-критерию ($p \geq 0{,}05$). \textit{Метрики снимаются на test-выборке RAVDESS из $n=96$ сэмплов; F1 оценивается дополнительно двумя независимыми внешними классификаторами выражений лица.}

\textbf{Гипотеза H2 (дополнительная).} Включение той же функции потерь в процедуру файнтюнинга SadTalker приводит к улучшению F1-меры эмоциональной согласованности на тестовой выборке не менее чем на 8~\% при сохранении ограничения на качество движений головы и общую визуальную целостность кадра, что демонстрирует методологическую воспроизводимость подхода на второй опорной архитектуре, отличающейся от Wav2Lip по типу представления (трёхмерная геометрия лица вместо двумерной маски). \textit{Критерии подтверждения:} $\Delta\text{F1} \geq +8$~\% ($p < 0{,}05$, критерий Макнемара) на тестовом подмножестве, оцениваемом независимыми внешними классификаторами по рендеренным видео.

\subsection{Методы исследования}

В работе применяются методы глубокого обучения (предобученные модели wav2vec~2.0 в варианте \texttt{superb/wav2vec2-base-superb-er} и TimeSformer в варианте \texttt{facebook/timesformer-base-finetuned-k400} для извлечения эмоциональных представлений из аудио- и видеомодальностей), а также методы кросс-модальной регуляризации, оптимальная конфигурация которой эмпирически устанавливается в пятиконфигурационной абляции. Для проверки статистической значимости применяется критерий Макнемара для категориальных метрик и парный $t$-критерий для метрик с количественной шкалой. Программная реализация выполнена на языке Python с применением библиотек PyTorch и Hugging Face Transformers; логирование экспериментов проведено в системе Weights \& Biases. Для снижения риска цикличности оценки применяется независимая внешняя оценка двумя классификаторами эмоций, не задействованными на этапах обучения и отбора моделей; оба внешних классификатора отобраны по результатам скрининга пяти кандидатов на тестовой выборке RAVDESS.

\subsection{Научная новизна}

Научная новизна исследования формулируется в границах проведённого эксперимента и состоит в следующем. Впервые показано, что повышение эмоциональной согласованности генерации говорящих лиц на четырёх эмоциях RAVDESS достижимо посредством plug-in функции потерь, применяемой в режиме файнтюнинга к опорной архитектуре Wav2Lip без модификации backbone и без явных эмоциональных меток на этапе инференса; обоснованность выбора компонент функции потерь (победила комбинация перекрёстной энтропии и KL-дистилляции от аудиоэнкодера к видеоэнкодеру) подтверждена пятиконфигурационным абляционным анализом, в котором наивная косинусная формулировка показала отрицательную $\Delta\text{F1}$ относительно baseline. Эффект подтверждён двумя независимыми внешними классификаторами, не задействованными в обучающей сигнальной цепи и предварительно отобранными в процедуре скрининга. В отличие от EmoTalk, EmoTalker, EmotiveTalk и EAT, использующих специализированные архитектуры или диффузионные стеки и оценивающих результат на тех же моделях, что использованы в loss, предложенный подход совместим с уже развёрнутыми моделями и реализует протокол внешней оценки, систематически отсутствующий в перечисленных работах. Распространение подхода на SadTalker как plug-in надстройки в текущем эксперименте не подтверждено; пределы переноса методологии на иные архитектуры обсуждаются в разделе перспектив.

\subsection{Теоретическая и практическая значимость}

\textbf{Теоретическая значимость} заключается в формализации задачи повышения эмоциональной согласованности при файнтюнинге как задачи кросс-модального контрастивного выравнивания на фиксированных энкодерах и в постановке протокола независимой внешней валидации, отделяющего вклад функции потерь от свойств обучающих метрик.

\textbf{Практическая значимость} определяется применимостью разработанной функции потерь к двум распространённым опорным архитектурам без переобучения с нуля, что снижает вычислительные затраты на адаптацию существующих систем и обеспечивает совместимость с уже развёрнутыми пайплайнами генерации говорящих лиц.

\subsection{Апробация результатов исследования}

Основные результаты диссертационного исследования представлены на 68-й Всероссийской научной конференции МФТИ (секция когнитивных технологий, 4 апреля 2026 г., Москва); работа принята для публикации в сборнике тезисов конференции (\url{https://conf.mipt.ru/conference/23659}). Программная реализация опубликована в открытом репозитории GitHub (\url{https://github.com/Katrin-Pochtar/The-Uncanny-Valley}); журналы экспериментов доступны в системе Weights \& Biases.

\subsection{Структура работы}

Диссертация состоит из введения, теоретического анализа, методологии исследования, главы результатов, перспектив и заключения, а также списка использованных источников и приложений. В теоретическом анализе систематизируются существующие подходы к синтезу говорящих лиц и формулируется теоретический пробел. В методологической главе описаны файнтюнинг аудио- и видеоэнкодеров, формализация и реализация кросс-модальной функции потерь, а также процедуры файнтюнинга Wav2Lip и SadTalker. В главе результатов представлены количественные результаты по гипотезам H1 и H2, абляционный анализ компонент функции потерь и независимая внешняя оценка. В разделе перспектив обсуждаются ограничения исследования и направления дальнейших работ. Заключение содержит обобщённые выводы по результатам проверки выдвинутых гипотез.

\section{Обозначения и сокращения}

\renewcommand{\arraystretch}{1.2}
\renewcommand{\tabcolsep}{0.3cm}
\begin{longtable}{|p{3.5cm}|p{12cm}|}
\hline
\textbf{Сокращение} & \textbf{Расшифровка} \\
\hline
\endfirsthead
\hline
\textbf{Сокращение} & \textbf{Расшифровка} \\
\hline
\endhead
ВКР & выпускная квалификационная работа \\
\hline
3DMM & трёхмерная морфируемая модель лица (3D Morphable Model) \\
\hline
AU & элементарная единица мимики (Action Unit) \\
\hline
CE & перекрёстная энтропия (Cross Entropy) \\
\hline
CNN & свёрточная нейронная сеть (Convolutional Neural Network) \\
\hline
FER & распознавание выражений лица (Facial Expression Recognition) \\
\hline
F1 & F1-мера, среднее гармоническое точности и полноты \\
\hline
GAN & генеративно-состязательная сеть (Generative Adversarial Network) \\
\hline
HuBERT & Hidden-Unit BERT - аудиоэнкодер с самообучением \\
\hline
KL & расхождение Кульбака - Лейблера (Kullback - Leibler divergence) \\
\hline
L1 & L1-норма ошибки реконструкции \\
\hline
LSE-C, LSE-D & Lip Sync Error - Confidence/Distance, метрики синхронизации губ \\
\hline
MEAD & Multi-view Emotional Audio-visual Dataset \\
\hline
MSP-DIM & модель размерной оценки эмоции (arousal/dominance/valence) \\
\hline
RAVDESS & Ryerson Audio-Visual Database of Emotional Speech and Song \\
\hline
SupCon & Supervised Contrastive Loss (Khosla et al., 2020) \\
\hline
SyncNet & модель оценки качества синхронизации губ \\
\hline
TimeSformer & пространственно-временной трансформер для видео \\
\hline
ViT & Vision Transformer \\
\hline
VAD & размерное представление эмоции: валентность/возбуждение/доминирование \\
\hline
VAE / VQ-VAE & вариационный автокодировщик / его векторно-квантованный вариант \\
\hline
Wav2Lip & модель аудио-управляемой синхронизации губ \\
\hline
Wav2Vec2 & самообучающаяся аудио-модель \\
\hline
W\&B & сервис логирования экспериментов Weights \& Biases \\
\hline
$\Delta$F1 & изменение F1-меры относительно базовой модели \\
\hline
$\Delta$LSE-C & изменение метрики LSE-C относительно базовой модели \\
\hline
H1, H2 & основная и дополнительная гипотезы исследования \\
\hline
\end{longtable}

\section{Теоретический анализ}

\subsection{Логика построения обзора литературы}

Настоящий раздел реализует системный обзор научной литературы по тематике аудио-управляемого синтеза говорящих лиц с акцентом на проблему эмоциональной согласованности. Структура обзора подчинена логике последовательного сужения от широкого контекста к узкой исследовательской постановке. На первом этапе рассматриваются психологические основы классификации эмоций, определяющие выбор представления эмоционального пространства. На втором этапе анализируются основные архитектурные классы существующих генеративных моделей и характер используемых ими функций потерь. На третьем этапе систематизируются подходы к интеграции эмоционального компонента в задачу синтеза говорящих лиц с разделением на категории по типу управляющего сигнала и степени модификации базовой архитектуры. На четвёртом этапе анализируются предобученные энкодеры эмоциональных представлений для аудио- и видеомодальностей. На пятом этапе рассматриваются датасеты, используемые в задачах распознавания эмоций и эмоционально-обусловленного синтеза. На шестом этапе рассматриваются методы кросс-модального контрастивного обучения, составляющие теоретическую основу предлагаемой функции потерь. На седьмом этапе анализируются метрики оценки качества генерации и эмоциональной согласованности с особым вниманием к методологическому риску цикличности оценки. На заключительном этапе проводится систематизация выявленных подходов в форме сравнительной таблицы и формулируется теоретический пробел, обосновывающий направление настоящего исследования.

\subsection{Психологические основы классификации эмоций}

Выбор представления эмоционального пространства в задачах распознавания и синтеза эмоционально-окрашенного контента опирается на классические психологические модели эмоций.

\textbf{Категориальная модель базовых эмоций.} Согласно теории, формализованной в работах [Ekman, 1992; Ekman, Friesen, 1978], существует ограниченный набор универсальных базовых эмоций, имеющих кросс-культурное проявление и характерные мимические корреляты. В классической формулировке выделяются шесть базовых эмоций (радость, печаль, гнев, страх, отвращение, удивление); расширенный вариант включает дополнительно нейтральное и презрительное состояния. Категориальная модель составляет основу большинства корпусов с эмоциональной разметкой (RAVDESS, CREMA-D, AffectNet) и применяется в качестве стандарта для оценки моделей распознавания и генерации эмоций.

\textbf{Размерная модель.} Альтернативный подход, восходящий к работам [Russell, 1980] и развитый в модели VAD (valence-arousal-dominance), представляет эмоциональное состояние не дискретной меткой, а точкой в непрерывном многомерном пространстве. Валентность характеризует приятность переживания, возбуждение - интенсивность активации, доминирование - степень контроля. Размерная модель обеспечивает более тонкую градацию состояний и применяется в корпусах MSP-Podcast и SEMAINE; в задачах распознавания речевых эмоций реализована в моделях [Wagner et al., 2023]. Преимуществом размерной модели является непрерывная природа представления, недостатком - менее интуитивная интерпретация и сложность построения корпусов с надёжной разметкой.

\textbf{Колесо эмоций Плутчика.} Модель [Plutchik, 1980] вводит восемь базовых эмоций, попарно противоположных, и допускает построение производных смешанных состояний. Указанная модель используется преимущественно в задачах анализа текста и менее распространена в аудиовизуальном анализе.

\textbf{Выбор представления для настоящего исследования.} В работе используется категориальная модель базовых эмоций в варианте, ограниченном четырьмя классами датасета RAVDESS: радость (happy), печаль (sad), гнев (angry) и отвращение (disgust). Сужение исходного восьмиклассового пространства RAVDESS до четырёх категорий мотивировано двумя соображениями. Во-первых, состояния <<нейтральное>> (neutral) и <<спокойное>> (calm) обладают слабой эмоциональной выраженностью и низкой различимостью при автоматической оценке мимики, что делает их включение методологически проблемным в задаче, где целевой эффект - именно повышение эмоционально окрашенной выразительности. Во-вторых, классы <<страх>> (fearful) и <<удивление>> (surprised) в студийной актёрской записи RAVDESS систематически конфузируются с радостью на уровне внешних предобученных классификаторов FER-производного происхождения, использование которых составляет ядро протокола внешней оценки в настоящей работе; включение указанных классов привело бы к интерференции между источниками шума оценивания и истинным эффектом функции потерь. Используемый набор из четырёх эмоций согласуется с распространённой практикой работ по эмоциональному синтезу [Ji et al., 2022; Wang et al., 2023] и обеспечивает интерпретируемое сопоставление матриц ошибок между внутренними и внешними классификаторами. Ограничения, связанные с указанным сужением эмоционального пространства, явно отражены в разделе, посвящённом границам исследования.

\subsection{Архитектурные классы моделей генерации говорящих лиц}

Современные модели аудио-управляемого синтеза говорящих лиц могут быть разделены на четыре архитектурных класса по способу представления и обновления лицевой геометрии: двумерные модели (image-space), трёхмерные модели на основе морфируемых представлений лица (3DMM-based), модели на основе нейронных радиационных полей и трёхмерного гауссова сплаттинга (NeRF / 3D Gaussian Splatting) и диффузионные модели.

\textbf{Двумерные image-space модели} оперируют непосредственно в пространстве пикселей входного кадра. Архитектура Wav2Lip [Prajwal et al., 2020] основана на парадигме <<кодировщик - декодер>> с дополнительным дискриминатором синхронизации губ (SyncNet), что обеспечивает точность синхронизации артикуляции свыше 97~\% на стандартных тестовых корпусах при работе с произвольным лицом. Модель оптимизирует функцию реконструкции L1 в нижней половине лица и совместно обучается с предобученным lipsync-экспертом, замороженным в процессе тренировки. Дальнейшие модификации, включая Wav2Lip-HQ [Mallikarjuna et al., 2025] и VividWav2Lip [Liu et al., 2024], направлены на повышение визуального качества посредством сегментации лица и применения сверхразрешающих сетей; явных механизмов управления эмоциональной выразительностью в указанных модификациях не вводится. К данному классу также относится модель MakeItTalk [Zhou et al., 2020], реализующая управление позой головы через предсказание ключевых точек лица в качестве промежуточного представления.

\textbf{Трёхмерные модели на основе морфируемой геометрии} разделяют процесс синтеза на два этапа: на первом из аудиосигнала предсказываются параметры трёхмерной геометрии лица (выражение, поза, освещение), на втором рендерятся видеокадры. Модель SadTalker [Zhang et al., 2023] использует двухветвевой подход: блок PoseVAE предсказывает движение головы, блок ExpNet - параметры выражений, после чего применяется промежуточная 3DMM-репрезентация. Подобная декомпозиция обеспечивает более естественные движения головы по сравнению с image-space моделями, однако эмоциональное содержание контролируется только опосредованно через параметры выражений, без явного связывания с эмоциональной семантикой аудиосигнала. Сходный двухэтапный подход применён в моделях Audio2Head [Wang et al., 2021], EAMM [Ji et al., 2022] и DiffPoseTalk [Sun et al., 2024]. Общим ограничением класса является зависимость качества выходного видео от точности промежуточной геометрической репрезентации.

\textbf{Модели на основе нейронных радиационных полей} реализуют синтез фотореалистичного видео через объёмное представление сцены, обусловленное аудиопризнаками. Семейство AD-NeRF [Guo et al., 2021], RAD-NeRF [Tang et al., 2022] и GeneFace [Ye et al., 2023] требует индивидуального обучения NeRF на видеоматериале конкретного субъекта, что ограничивает обобщение на произвольное лицо, но обеспечивает высокую визуальную точность для целевого субъекта. Подходы на основе 3D Gaussian Splatting (InsTaG [Li et al., 2025], EmoDiffTalk [Liu et al., 2026]) сокращают время обучения и улучшают параметрическую управляемость, в том числе через интеграцию управления Action Units. Эмоциональный контроль в данном классе моделей реализуется преимущественно через явные текстовые или AU-промпты на этапе генерации, что требует наличия дополнительной разметки и не соответствует постановке задачи извлечения эмоции непосредственно из аудиосигнала.

\textbf{Диффузионные модели} реализуют синтез как последовательный процесс денойзинга, обусловленный аудиопризнаками. Модель DiffTalk [Shen et al., 2023] применяет латентное диффузионное моделирование для портретной анимации; EMO [Tian et al., 2024] комбинирует диффузионный процесс с временной свёрткой для генерации длительных эмоционально насыщенных последовательностей; EmotiveTalk [Wang et al., 2025] вводит декуплирование аудиопризнаков для раздельного контроля артикуляции и эмоций. Диффузионные подходы достигают высокого качества генерации, однако требуют существенных вычислительных ресурсов и, как правило, замены всего генеративного стека, что ограничивает их применимость к задаче файнтюнинга уже развёрнутых моделей.

Систематизация показывает, что несмотря на архитектурное разнообразие, перечисленные классы моделей либо не предусматривают явного управления эмоциональной выразительностью (классы 1 и 2), либо требуют замены архитектуры или явных эмоциональных промптов на этапе инференса (классы 3 и 4). Опорные архитектуры Wav2Lip и SadTalker, выбранные в качестве объектов настоящего исследования, относятся к первым двум классам и представляют собой широко применяемые в практике решения, для которых задача надстройки эмоционального контроля без модификации архитектуры представляет самостоятельный исследовательский интерес.

\subsection{Эмоциональное измерение в синтезе говорящих лиц}

Подходы к интеграции эмоционального компонента в задачу синтеза говорящих лиц могут быть классифицированы по типу управляющего сигнала и по степени модификации базовой архитектуры. Указанная классификация позволяет позиционировать настоящее исследование относительно существующих работ.

\textbf{Управление через явные категориальные метки} реализовано в семействе моделей, обученных на корпусах с эмоциональной разметкой. Датасет MEAD [Wang et al., 2020] (Multi-view Emotional Audio-visual Dataset) содержит видеозаписи 60 актёров с восемью базовыми эмоциями и тремя уровнями интенсивности; на его основе разработаны специализированные эмоциональные модели, включая EAMM [Ji et al., 2022] и EAT [Wang et al., 2023]. Модель EAT (Efficient Emotional Adaptation for Audio-Driven Talking-Head Generation) реализует адаптацию через дополнительные эмоциональные модули, требующие явной метки эмоции на этапе инференса. Подобный подход обеспечивает прямой контроль выразительности, но требует дополнительной разметки в применении и не масштабируется на сценарии, где эмоция должна выводиться непосредственно из аудиосигнала.

\textbf{Управление через специализированную архитектуру} реализовано в моделях, проектируемых с эмоциональным компонентом как часть целевой функции обучения. EmoTalk [Peng et al., 2023] выполняет разделение эмоциональных и лингвистических признаков для трёхмерной анимации лица посредством специальных энкодеров и обучается с нуля. EmoCAST [Jiang et al., 2025] применяет диффузионный фреймворк с управлением через текстовые описания эмоций. EMOCA [Daněček et al., 2022] решает обратную задачу - реконструкцию трёхмерной геометрии лица из видео с учётом эмоциональной информации - и не предназначена для аудио-управляемого синтеза. Указанные подходы либо требуют полной переработки архитектуры, либо относятся к иной постановке задачи.

\textbf{Управление через Action Units} опирается на систему кодирования лицевых действий (Facial Action Coding System, FACS), формализованную в работе [Ekman, Friesen, 1978]. Action Units представляют собой элементарные мимические движения, имеющие нейроанатомическую интерпретацию и обеспечивающие компактное структурированное описание выражения лица. Модели AUHead [Lyu et al., 2026], MF-ETalk [Wang et al., 2025] и EmoDiffTalk [Liu et al., 2026] используют AU как промежуточное представление между аудиосигналом и видеовыходом. AU-предсказательные модели обеспечивают интерпретируемость и анатомическую корректность, однако требуют вычислительно дорогих этапов промежуточной разметки и систем предсказания AU из аудио.

\textbf{Управление через кросс-модальное выравнивание без явных меток на инференсе} представлено меньшим числом работ. EmoTalker [Shen et al., 2024] реализует кросс-модальный трансформер, обученный с нуля, который извлекает эмоциональную информацию непосредственно из аудио и применяет её при генерации без необходимости в явной метке. CEM-Net [Wu et al., 2025] решает связанную задачу разрешения конфликта между референсным изображением и аудиосигналом посредством сети памяти. Подобные подходы концептуально близки к настоящему исследованию, однако предполагают обучение специализированной архитектуры с нуля, а не файнтюнинг существующей опорной модели.

\textbf{Полемика в литературе.} В исследовательском сообществе наблюдаются различные позиции относительно компромисса между интерпретируемостью и обобщающей способностью. Сторонники AU-управления [Lyu et al., 2026; Liu et al., 2026] аргументируют необходимость явного структурированного представления для обеспечения контроля и анатомической корректности, опираясь на нейропсихологическую обоснованность системы FACS. Сторонники end-to-end кросс-модальных подходов [Shen et al., 2024; Wang et al., 2025] указывают на ограниченность дискретных категориальных или AU-представлений для передачи тонких эмоциональных оттенков и отстаивают применение непрерывных аудио-производных представлений. Перцепционные исследования [Suma et al., 2023; Zhou et al., 2025] подтверждают значимость согласованности голосового и мимического каналов независимо от конкретного механизма управления, опираясь на классическое наблюдение мультимодальной интеграции в восприятии речи [McGurk, MacDonald, 1976].

Систематизация выявляет отсутствие в литературе работ, в которых эмоциональная согласованность достигается посредством надстройки над уже существующими опорными архитектурами без их замены и без явной разметки эмоций на этапе инференса. Указанная ниша определяет постановку настоящего исследования.

\subsection{Аудио- и видеоэнкодеры эмоциональных представлений}

Извлечение эмоциональных представлений из соответствующих модальностей реализуется специализированными предобученными моделями, обучение которых на корпусах с эмоциональной разметкой обеспечивает высокую разделяющую способность в пространстве эмбеддингов.

\textbf{Аудиоэнкодеры.} Архитектуры на основе самообучения с маскировкой (self-supervised learning) показали высокую эффективность в задачах распознавания речевых эмоций. Модель wav2vec~2.0 [Baevski et al., 2020] обучается на сырых аудиоволнах через контрастивную задачу предсказания квантованных представлений; HuBERT [Hsu et al., 2021] заменяет контрастивную задачу на маскированное предсказание кластерных меток, полученных кластеризацией k-means на признаках MFCC. После файнтюнинга на эмоциональных корпусах (RAVDESS, IEMOCAP) указанные модели достигают F1-меры 0,75-0,90 на классификации базовых эмоций [Pepino et al., 2021]. Расширение AV-HuBERT [Shi et al., 2022] обобщает HuBERT на аудиовизуальную модальность через совместное маскированное обучение и применяется в задачах распознавания речи в шумных условиях.

\textbf{Видеоэнкодеры.} Архитектура Vision Transformer [Dosovitskiy et al., 2021] обобщает механизм самовнимания на изображения через разбиение на патчи и линейное проецирование. TimeSformer [Bertasius et al., 2021] адаптирует ViT к видео через факторизацию пространственного и временного внимания, что снижает вычислительную сложность и обеспечивает применимость на длинных последовательностях кадров. Для распознавания выражений лица применяются специализированные ViT, предобученные на корпусах AffectNet [Mollahosseini et al., 2017] и FER+ [Barsoum et al., 2016]. Альтернативный класс представлен моделями MARLIN [Cai et al., 2023] и Video Masked Autoencoders [Tong et al., 2022], реализующими маскированное автокодирование для видеомодальности. Для задач, требующих интерпретируемых выражений, применяются также модели предсказания AU, обученные на корпусах с FACS-разметкой.

\textbf{Кросс-модальные представления.} CLIP [Radford et al., 2021] продемонстрировал эффективность контрастивного обучения на парах <<изображение - текст>> для построения общего семантического пространства. ALIGN [Jia et al., 2021] масштабирует подход на корпусе из 1,8 миллиарда пар. CAV-MAE [Gong et al., 2023] адаптирует масочное автокодирование к аудиовизуальной модальности и реализует совместное обучение через комбинацию реконструкции и контрастивного выравнивания. Применение принципа CLIP к паре <<эмоциональное аудио - эмоциональное видео>> при синтезе говорящих лиц в литературе систематически не исследовано; имеющиеся работы в области аудио-визуального согласования (например, [Owens, Efros, 2018]) сфокусированы на задачах распознавания речи и идентификации говорящего, а не на эмоциональной семантике в задаче генерации.

\textbf{Выбор энкодеров для настоящего исследования.} В качестве отправной точки рассмотрены три семейства предобученных аудиомоделей (Wav2Vec2 \texttt{base}, Wav2Vec2 \texttt{large}, HuBERT \texttt{large}, а также SUPERB-варианты, дообученные на корпусе IEMOCAP) и одно семейство видеомоделей (TimeSformer, варианты с числом кадров 4/8/12/16 и различными темпами обучения). По итогам сравнительной оценки на валидационной выборке RAVDESS (4-классовая постановка, разбиение по актёрам) для дальнейшей работы отобраны: \texttt{superb/wav2vec2-base-superb-er} в качестве аудиоэнкодера (валидационная макро-F1 0,826) и \texttt{facebook/timesformer-base-finetuned-k400} с 8-кадровой выборкой в качестве видеоэнкодера (валидационная макро-F1 0,790). На независимой тестовой выборке указанные модели достигают макро-F1 0,739 (аудио) и 0,876 (видео), что фиксирует <<потолок>> эмоциональной разделимости, доступной любой модели генерации, оцениваемой через данные энкодеры. Подробности гиперпараметрического перебора и полная таблица конфигураций приведены в разделе результатов и в ноутбуке [02\_train\_encoders\_4emotions.ipynb](02\_train\_encoders\_4emotions.ipynb). Выходы классификационных голов указанных моделей (4-мерные логиты эмоций) используются как целевые сигналы в формулируемой далее кросс-модальной функции потерь. Для оценки на этапе финального тестирования применяются независимые внешние классификаторы выражений лица, не задействованные ни на одном этапе обучения и отбора моделей; обоснование выбора и процедура скрининга приведены в подразделе, посвящённом метрикам.

\subsection{Датасеты для эмоционально-обусловленного синтеза}

Качество эмоционально-обусловленных моделей генерации существенно зависит от характеристик обучающих корпусов: количества эмоциональных категорий, числа и разнообразия субъектов, условий записи и наличия парной аудиовизуальной разметки. В литературе сложился ограниченный круг датасетов, регулярно применяемых в задачах распознавания и синтеза эмоций.

\textbf{RAVDESS (Ryerson Audio-Visual Database of Emotional Speech and Song)} [Livingstone, Russo, 2018] содержит 2880 аудиовизуальных записей 24 актёров (12 мужчин и 12 женщин), произносящих две стандартные фразы в восьми эмоциональных состояниях (нейтральное, спокойное, радость, печаль, гнев, страх, отвращение, удивление) с двумя уровнями интенсивности. Запись осуществлялась в студийных условиях с фиксированной позой головы и контролируемым освещением, что обеспечивает стандартизацию, но ограничивает внешнюю валидность относительно условий in-the-wild. RAVDESS является наиболее широко используемым сбалансированным корпусом для аудиовизуального распознавания эмоций с парной разметкой речи и видео и применяется в качестве основного датасета настоящего исследования.

\textbf{MEAD (Multi-view Emotional Audio-visual Dataset)} [Wang et al., 2020] содержит видеозаписи 60 актёров с восемью базовыми эмоциями и тремя уровнями интенсивности, снятые с семи ракурсов. MEAD превосходит RAVDESS по объёму и разнообразию субъектов и применяется в специализированных эмоциональных моделях [Ji et al., 2022; Wang et al., 2023]. Использование MEAD в настоящем исследовании выходит за пределы экспериментально подтверждённого охвата и составляет одно из направлений дальнейших работ.

\textbf{CREMA-D (Crowd-sourced Emotional Multimodal Actors Dataset)} [Cao et al., 2014] содержит 7442 записи 91 актёра разнородного этнокультурного происхождения с шестью базовыми эмоциями и четырьмя уровнями интенсивности. По сравнению с RAVDESS обеспечивает большее разнообразие субъектов; внешняя валидация на CREMA-D также рассматривается как направление дальнейших исследований.

\textbf{IEMOCAP (Interactive Emotional Dyadic Motion Capture)} [Busso et al., 2008] содержит около 12 часов аудиовизуальных диалогов с двойной разметкой (категориальной и размерной), полученной от нескольких аннотаторов. IEMOCAP применяется преимущественно в задачах распознавания речевых эмоций и менее подходит для задач синтеза говорящих лиц вследствие диалогового характера записи и отсутствия фиксированных стандартизованных фраз.

\textbf{AffectNet} [Mollahosseini et al., 2017] и \textbf{FER+} [Barsoum et al., 2016] являются крупномасштабными корпусами изображений лиц с эмоциональной разметкой (около 1 миллиона и 35 тысяч изображений соответственно). Указанные корпуса используются для предобучения видеоэнкодеров и независимых классификаторов выражений лица; в настоящем исследовании внешний классификатор основан на FER-производном корпусе, что обеспечивает независимость от RAVDESS.

\textbf{Обоснование выбора RAVDESS.} Выбор RAVDESS в качестве основного датасета мотивирован следующим: (1) сбалансированное представление восьми категорий эмоций при одинаковом числе записей на категорию обеспечивает корректное обучение и оценку без необходимости weighted sampling; (2) парная аудиовизуальная разметка с фиксированными фразами устраняет конфаундинг между лингвистическим содержанием и эмоциональной окраской; (3) контролируемые условия записи минимизируют влияние освещения и позы головы на извлекаемые эмоциональные представления; (4) корпус является открыто доступным и широко используемым в литературе, что обеспечивает воспроизводимость и сопоставимость результатов. Ограничения, связанные со студийным характером записи и относительно небольшим числом субъектов, явно отражены в разделе, посвящённом границам исследования.

\subsection{Контрастивное обучение и кросс-модальное выравнивание}

Теоретическую основу предлагаемой функции потерь составляют методы контрастивного обучения, цель которых - построение представлений, в которых семантически близкие объекты сближаются, а семантически различные - раздвигаются в общем векторном пространстве.

\textbf{Фундаментальные формулировки.} Целевая функция InfoNCE [van den Oord et al., 2018] обеспечивает построение representation через максимизацию нижней границы взаимной информации между связанными представлениями посредством классификации позитивной пары среди множества негативных. Метод SimCLR [Chen et al., 2020] продемонстрировал применимость подхода к самообучению на изображениях, используя augmentations как источник позитивных пар. Метод MoCo [He et al., 2020] вводит momentum encoder и dynamic queue для увеличения числа доступных негативных примеров, что особенно существенно при ограниченном размере батча.

\textbf{Supervised contrastive learning.} Метод SupCon [Khosla et al., 2020] расширяет контрастивную постановку на supervised режим: позитивами для якоря $i$ являются все объекты с той же меткой класса, а не только augmentations того же объекта. Подобная формулировка обеспечивает построение более структурированного embedding-пространства и в ряде задач классификации показала превосходство над cross-entropy. Применение SupCon к кросс-модальной постановке (<<аудио - видео>>) с использованием эмоциональных меток как источника позитивных пар реализуется в настоящем исследовании в качестве основной формулировки функции потерь и сопоставляется в абляции с альтернативными вариантами.

\textbf{Кросс-модальное выравнивание.} CLIP [Radford et al., 2021] обучает совместное представление изображений и текста через симметричную InfoNCE-постановку с использованием в качестве позитивных пар связанных по природе объектов и в качестве негативных - несвязанных пар в батче. Применительно к аудио-видео модальности AV-HuBERT [Shi et al., 2022] и CAV-MAE [Gong et al., 2023] реализуют выравнивание через совместное маскированное обучение, однако не специализированы на эмоциональной семантике; используемые ими позитивные пары формируются по принципу синхронности во времени, а не по эмоциональной общности.

\textbf{Риск коллапса проекций.} При попытке наивного применения косинусного сближения к парам <<аудио - видео>> одного и того же сэмпла без использования негативных примеров возникает риск тривиального коллапса: проекционные слои сходятся к константному отображению, обеспечивающему максимум косинусного сходства без передачи семантической информации [Chen, He, 2021]. Механизмы предотвращения коллапса в литературе по самообучению включают применение остановленного градиента (stop-gradient), использование предсказателя на одной из ветвей (BYOL [Grill et al., 2020]) и явное введение негативных пар. В настоящем исследовании предотвращение коллапса достигается посредством перехода от наивной косинусной формулировки к контрастивной с явными негативными парами, формируемыми на основе несовпадения эмоциональных меток. Указанный методологический выбор подтверждён эмпирически в рамках абляционного анализа.

\subsection{Метрики оценки качества генерации и эмоциональной согласованности}

Оценка моделей синтеза говорящих лиц традиционно опирается на несколько групп метрик, характеризующих различные аспекты качества генерации.

\textbf{Метрики синхронизации губ.} Стандарт де-факто составляют метрики LSE-C (Lip-Sync Error - Confidence) и LSE-D (Lip-Sync Error - Distance), вычисляемые с применением модели SyncNet [Chung, Zisserman, 2016]. SyncNet представляет собой двухветвевую сеть, обученную на аудиовизуальных парах для предсказания временного смещения; LSE-C измеряет уверенность синхронизации, LSE-D - нормированное расстояние в общем embedding-пространстве. Метрики SyncNet используются в качестве первичного бенчмарка в работах [Prajwal et al., 2020; Zhang et al., 2023; Ye et al., 2023] и составляют ограничение на допустимое ухудшение качества артикуляции, формализованное в гипотезах H1 и H2 настоящего исследования.

\textbf{Метрики визуального качества.} Frechét Inception Distance [Heusel et al., 2017] оценивает расстояние между распределениями признаков сгенерированных и эталонных изображений в пространстве предобученной InceptionV3. Frechét Video Distance [Unterthiner et al., 2018] обобщает FID на видео через применение Inflated 3D ConvNet. LPIPS [Zhang et al., 2018] оценивает перцепционное сходство в пространстве признаков предобученной нейронной сети с обучением калибровочных весов на парных данных с человеческими оценками. Указанные метрики характеризуют общее визуальное качество, но не специфичны к эмоциональной выразительности и не позволяют оценить корректность эмоциональной семантики сгенерированного видео.

\textbf{Метрики эмоциональной согласованности.} Стандартный подход состоит в применении предобученного классификатора эмоций к сгенерированному видео и сопоставлении предсказанной метки с целевой эмоцией; в качестве агрегирующей метрики используется F1-мера или accuracy [Ji et al., 2022; Wang et al., 2023]. В отличие от метрик синхронизации губ, для оценки эмоциональной согласованности в литературе не сложилось универсального бенчмарка, и конкретный выбор классификатора (его обучающие данные, архитектура, версия) существенно влияет на результат, что затрудняет прямое сравнение между работами.

\textbf{Проблема цикличности оценки.} Если эмоциональный классификатор, используемый в качестве компонента функции потерь при обучении (например, для извлечения целевых эмоциональных представлений), одновременно применяется и для финальной оценки эмоциональной согласованности сгенерированного видео, возникает риск кругового рассуждения: модель оптимизируется под распределение, которое тот же классификатор успешно распознаёт, что не гарантирует переноса на независимое восприятие. Указанное методологическое ограничение систематически отмечается в обзорной литературе [Jiang et al., 2024], однако в практике многих работ по эмоциональному синтезу мер по его явному устранению не предусмотрено. В настоящем исследовании в качестве методологического ответа на данный риск проведён скрининг шести предобученных внешних классификаторов выражений лица (\texttt{dima806/facial\_emotions\_image\_detection}, \texttt{trpakov/vit-face-expression}, \texttt{motheecreator/vit-Facial-Expression-Recognition}, \texttt{Rajaram1996/FacialEmoRecog}, \texttt{RickyIG/emotion\_face\_image\_classification}, \texttt{prithivMLmods/Facial-Emotion-Detection-SigLIP2}); последний исключён из-за несовместимости меток (отсутствие класса \textit{disgust}). По результатам скрининга на тестовой выборке RAVDESS (см. раздел результатов) для финальной оценки отобраны два классификатора с максимальной макро-F1 на четырёх целевых эмоциях - \texttt{motheecreator/vit-Facial-Expression-Recognition} и \texttt{Rajaram1996/FacialEmoRecog}; оба обучены на FER-производных корпусах, не пересекающихся с RAVDESS, что обеспечивает независимость оценочного сигнала. Перенос направления эффекта между двумя различающимися по архитектуре и обучающему корпусу внешними классификаторами фиксируется в виде отдельной диагностической метрики (cross-external agreement).

\textbf{Перцепционные метрики.} Слепое A/B тестирование и оценка по шкале Ликерта [Likert, 1932] остаются золотым стандартом для оценки субъективного качества генерации [Zhou et al., 2025]. Объёмная перцепционная оценка с привлечением группы респондентов выходит за пределы настоящего исследования и составляет одно из направлений дальнейших работ, обозначенных в разделе перспектив.

\subsection{Систематизация и формулировка теоретического пробела}

Проведённый обзор позволяет систематизировать существующие подходы к эмоционально-обусловленной генерации говорящих лиц по нескольким критериям, существенным для настоящего исследования: требуется ли замена базовой архитектуры (backbone); требуется ли явная метка эмоции на этапе инференса; специализирован ли подход на эмоциональной семантике; реализована ли независимая внешняя валидация. Результаты систематизации представлены в таблице~\ref{tab:approaches_comparison}.

\renewcommand{\arraystretch}{1.3}
\renewcommand{\tabcolsep}{0.15cm}
\begin{table}[H]
\caption{Сравнение подходов к эмоционально-обусловленной генерации говорящих лиц}
\label{tab:approaches_comparison}
\centering
\small
\begin{tabular}{|p{4.5cm}|c|c|c|c|}
\hline
\textbf{Подход} & \textbf{Замена backbone} & \textbf{Явные метки на инференсе} & \textbf{Эмоциональная специализация} & \textbf{Внешняя валидация} \\
\hline
Wav2Lip [Prajwal et al., 2020] & - & - & нет & - \\
\hline
SadTalker [Zhang et al., 2023] & - & - & нет & - \\
\hline
EmoTalk [Peng et al., 2023] & да & да & да & нет \\
\hline
EAT [Wang et al., 2023] & да & да & да & нет \\
\hline
EmoTalker [Shen et al., 2024] & да & нет & да & нет \\
\hline
EmotiveTalk [Wang et al., 2025] & да & нет & да & нет \\
\hline
EmoCAST [Jiang et al., 2025] & да & да (текст) & да & нет \\
\hline
AUHead [Lyu et al., 2026] & да & да (AU) & да & нет \\
\hline
EmoDiffTalk [Liu et al., 2026] & да & да (AU) & да & нет \\
\hline
EMOCA [Daněček et al., 2022] & - & - & да & нет \\
\hline
\textbf{Настоящее исследование} & \textbf{нет} & \textbf{нет} & \textbf{да} & \textbf{да} \\
\hline
\end{tabular}
\end{table}

Анализ сравнительной таблицы выявляет три обстоятельства, существенных для постановки исследования. Во-первых, ни один из рассмотренных подходов не сочетает свойство plug-in применимости к существующим опорным архитектурам с независимой внешней валидацией результата. Во-вторых, подходы, обеспечивающие эмоциональную согласованность без явных меток на этапе инференса, требуют замены базовой архитектуры, что несовместимо с задачей адаптации уже развёрнутых систем. В-третьих, систематический контроль риска цикличности оценки в литературе не реализован: даже работы, явно адресующие эмоциональную согласованность, оценивают результат теми же моделями, что использованы в обучающей сигнальной цепи.

На основании выявленных обстоятельств формулируется \textbf{теоретический пробел}: в литературе отсутствует функция потерь, реализуемая как plug-in надстройка над опорными архитектурами Wav2Lip и SadTalker и обеспечивающая повышение эмоциональной согласованности с верификацией результата независимым внешним классификатором, не задействованным в обучающей цепи.

Сужение пробела до plug-in подхода для двух конкретных опорных архитектур мотивировано как теоретически, так и практически. Теоретически такое сужение позволяет сформулировать опровержимую гипотезу с измеримыми порогами, что соответствует требованиям к научной корректности. Практически - уже развёрнутые системы цифровых аватаров используют именно опорные архитектуры рассматриваемых классов, и переход на специализированные эмоциональные модели (EmoTalk, EmotiveTalk и аналогичные) требует существенных архитектурных и инфраструктурных изменений. Снижение барьера адаптации существующих систем к задаче эмоционально-согласованного синтеза составляет практическую значимость настоящего исследования.

Сформулированный теоретический пробел естественным образом ведёт к гипотезам H1 (для опорной архитектуры Wav2Lip) и H2 (для опорной архитектуры SadTalker), приведённым во Введении. Ожидается, что предлагаемая plug-in функция потерь обеспечит статистически значимое улучшение F1-меры эмоциональной согласованности при сохранении LSE-C в пределах допустимого ограничения, причём указанный эффект подтвердится не только по внутренним метрикам, но и при оценке независимым внешним классификатором. Подтверждение или опровержение указанных гипотез составляет предмет эмпирической части настоящей работы.


\section{Методология и методика исследования}

\subsection{Общая структура методики}

Методика исследования организована как последовательность четырёх взаимосвязанных этапов, выстроенных от подготовки данных и опорных моделей через формализацию функции потерь к экспериментальной проверке гипотез на двух опорных архитектурах. Каждый этап реализован отдельным ноутбуком в открытом репозитории; нумерация ноутбуков соответствует порядку выполнения в общем пайплайне:
\begin{enumerate}
\item Подготовка данных и разбиение по актёрам - [01\_data\_preprocessing.ipynb](01\_data\_preprocessing.ipynb).
\item Файнтюнинг аудио- и видеоэнкодеров эмоций (опорные модели функции потерь) - [02\_train\_encoders\_4emotions.ipynb](02\_train\_encoders\_4emotions.ipynb), [03\_emotion\_utils.ipynb](03\_emotion\_utils.ipynb).
\item Абляционный анализ компонент функции потерь и установление <<потолка>> разделимости энкодеров - [04a\_wav2lip\_loss\_ablation.ipynb](04a\_wav2lip\_loss\_ablation.ipynb), [04b\_encoder\_ceiling.ipynb](04b\_encoder\_ceiling.ipynb).
\item Скрининг внешних классификаторов и подготовка протокола внешней оценки - [external\_screening.ipynb](external\_screening.ipynb), [05\_external\_evaluation.ipynb](05\_external\_evaluation.ipynb).
\item Файнтюнинг опорных архитектур с применением победившей в абляции функции потерь и проверка гипотез - [06\_finetune\_wav2lip.ipynb](06\_finetune\_wav2lip.ipynb), [07\_finetune\_sadtalker.ipynb](07\_finetune\_sadtalker.ipynb), [07a\_sadtalker\_loss\_ablation.ipynb](07a\_sadtalker\_loss\_ablation.ipynb).
\end{enumerate}

Эксперимент в строгом смысле составляет этапы 3 и 5: на этапе 3 непосредственно манипулируется состав функции потерь при фиксированной архитектуре и данных, на этапе 5 - весовой коэффициент функции потерь при фиксированной её формулировке и опорной архитектуре. Все остальные этапы являются подготовительными или валидационными.

\subsection{Датасет и протокол разбиения}

В качестве основного корпуса используется RAVDESS (Ryerson Audio-Visual Database of Emotional Speech and Song) [Livingstone, Russo, 2018], публично доступный по лицензии CC BY-NA-SC 4.0. Из исходных восьми эмоциональных классов отобраны четыре: \textit{happy}, \textit{sad}, \textit{angry}, \textit{disgust}; обоснование выбора приведено в теоретическом разделе. Все 24 актёра RAVDESS разбиты на три непересекающиеся группы по идентификаторам субъектов (\textit{actor-aware split}), что исключает утечку признаков говорящего между обучающей, валидационной и тестовой выборками. Итоговые объёмы: обучение - 544 сэмпла, валидация - 96, тестирование - 96 (по 24 на каждую из четырёх эмоций). Один сэмпл представляет собой пару <<синхронизованная аудиозапись~/~видеозапись из 25~fps>>, обработанная скриптами этапа~1: видео нормализовано к разрешению $96\times 96$ для совместимости с архитектурой Wav2Lip и к тензору 8 кадров для видеоэнкодера; аудио ресэмплировано к 16~kHz и спектрально преобразовано в мел-спектрограмму по гиперпараметрам, совместимым с предобученным lipsync-экспертом SyncNet, используемым в качестве LSE-C/LSE-D-метрик.

\subsection{Опорные энкодеры эмоций (этап 2)}

Аудиоэнкодер и видеоэнкодер обучаются раздельно - как 4-классовые классификаторы эмоций - с целью получения двух независимых каналов <<эмоционального сигнала>>, выходы которых используются на этапе 5 в качестве целевых распределений в кросс-модальной функции потерь.

\textit{Гиперпараметрический перебор.} Для аудиомодальности рассмотрено 11 конфигураций (варианты \texttt{wav2vec2-base}, \texttt{wav2vec2-large}, \texttt{hubert-large} и SUPERB-вариантов с разными темпами обучения, разными окнами и режимами заморозки/класс-балансировки); для видео - 11 конфигураций \texttt{timesformer-base-finetuned-k400} с числом кадров 4/8/12/16 и темпами обучения от $1\!\cdot\!10^{-5}$ до $5\!\cdot\!10^{-5}$. Все запуски - 30 эпох с ранней остановкой по валидационной F1, оптимизатор AdamW, scheduler CosineAnnealingLR, batch-size 8, label smoothing 0,1, weight decay 0,01, seed 42. Логирование - W\&B (\texttt{uncanny-valley-encoders-4emo}). Полная таблица результатов перебора приведена в разделе результатов.

\textit{Финальный выбор.} По итогам перебора отобраны:
\begin{itemize}
\item \textbf{Аудиоэнкодер:} \texttt{superb/wav2vec2-base-superb-er}, темп обучения $3\!\cdot\!10^{-5}$, длина окна по умолчанию (валидационная макро-F1 0,826).
\item \textbf{Видеоэнкодер:} \texttt{facebook/timesformer-base-finetuned-k400}, 8 кадров на сэмпл, темп обучения $3\!\cdot\!10^{-5}$ (валидационная макро-F1 0,790).
\end{itemize}

Обе модели на этапах 3 и 5 \textit{зафиксированы} (\texttt{requires\_grad=False}); их выходы (4-мерные логиты эмоций) используются исключительно как сигнал в функции потерь, без дальнейшего обновления весов.

\subsection{Формулировка кросс-модальной функции потерь}

Функция потерь, дополняющая стандартную реконструкцию L1 при файнтюнинге опорной архитектуры $\mathcal{G}$, имеет общий вид
\begin{equation}\label{eq:totalloss}
\mathcal{L} = \mathcal{L}_{\text{recon}} + s \cdot \big(\lambda_{\text{ce}}\,\mathcal{L}_{\text{CE}} + \lambda_{\cos}\,\mathcal{L}_{\cos} + \lambda_{\text{KL}}\,\mathcal{L}_{\text{KL}}\big),
\end{equation}
где
\begin{itemize}
\item $\mathcal{L}_{\text{recon}} = \|\hat{V} - V\|_1$ - стандартная L1-реконструкция нижней половины кадра, унаследованная от Wav2Lip;
\item $\mathcal{L}_{\text{CE}} = \mathrm{CE}(\phi_v(\hat{V}),\, y)$ - перекрёстная энтропия между логитами видеоэнкодера на сгенерированном кадре и истинной меткой эмоции $y$;
\item $\mathcal{L}_{\cos} = 1 - \cos\!\bigl(\phi_v(\hat V),\, \phi_a(A)\bigr)$ - косинусное сближение признаков сгенерированного видео и аудиосигнала;
\item $\mathcal{L}_{\text{KL}} = \tau^2\,\mathrm{KL}\!\bigl(\sigma(\phi_v(\hat V)/\tau)\,\|\,\sigma(\phi_a(A)/\tau)\bigr)$ - KL-дистилляция, переносящая мягкое распределение классов с аудиоканала на видеоканал, температура $\tau=2{,}0$;
\item $s$ - общий масштабный коэффициент, $\lambda_{*}$ - покомпонентные веса, выбираемые в абляции.
\end{itemize}

Выбор именно такой формы продиктован тремя соображениями. Во-первых, формулировка \eqref{eq:totalloss} остаётся \textit{plug-in}: она не требует изменений ни в архитектуре $\mathcal{G}$, ни в её прямом проходе. Во-вторых, набор компонент покрывает три качественно различных режима выравнивания: явная супервизия по меткам ($\mathcal{L}_{\text{CE}}$), бесмаркерное выравнивание представлений ($\mathcal{L}_{\cos}$) и непрерывный перенос распределения с одной модальности на другую ($\mathcal{L}_{\text{KL}}$). В-третьих, оптимальное сочетание компонент не известно a priori и устанавливается эмпирически на этапе абляции.

\subsection{Абляционный анализ функции потерь (этап 3, проверка корректности формулировки)}

Цель этапа - эмпирически выбрать состав функции потерь, который при минимальной деградации L1-реконструкции даёт наибольший прирост F1 на валидационной выборке. Поскольку речь о выборе формулы, а не порога, испытываются пять конфигураций:
\begin{enumerate}
\item \textbf{abl-baseline} - $\lambda_{\text{ce}}=\lambda_{\cos}=\lambda_{\text{KL}}=0$ (только реконструкция).
\item \textbf{abl-cos-only} - $\lambda_{\cos}=0{,}2$, остальные веса нулевые. Контроль на коллапс наивной косинусной формулировки в отсутствие негативных пар.
\item \textbf{abl-ce-only} - $\lambda_{\text{ce}}=0{,}05$.
\item \textbf{abl-ce-cos} - $\lambda_{\text{ce}}=0{,}05$, $\lambda_{\cos}=0{,}2$.
\item \textbf{abl-ce-kl} - $\lambda_{\text{ce}}=0{,}05$, $\lambda_{\text{KL}}=0{,}1$, $\tau=2{,}0$.
\end{enumerate}

Все конфигурации обучаются на одном и том же seed (42) на 70 эпохах с ранней остановкой по 8 эпохам без улучшения валидационной F1 (для эмоциональных конфигураций) либо валидационной L1 (для baseline). Опорная архитектура - Wav2Lip; все слои трансформера и SyncNet-эксперт заморожены, обновляются только параметры аудио- и видеоэнкодеров и декодера Wav2Lip. Полная техническая спецификация запусков и логи доступны в проекте \texttt{uncanny-valley-wav2lip-ablation} в W\&B.

Для каждой конфигурации регистрируются: динамика валидационной L1 и F1 по эпохам; финальное значение L1, F1, accuracy на валидационной выборке; статистическая значимость отличия от baseline (для F1 - критерий Макнемара; для L1 - парный $t$-критерий и критерий Уилкоксона). Победившая по F1 конфигурация переходит на этап 5 для проверки гипотез на тестовой выборке.

\subsection{Поиск весового коэффициента и проверка H1 (этап 5, Wav2Lip)}

Победившая в абляции комбинация (CE+KL) применяется к Wav2Lip с переменным масштабом $s\in\{0{,}1;\,0{,}25;\,0{,}4;\,0{,}5\}$. Для каждой точки сохраняется чекпойнт по максимуму валидационной F1. Правило отбора финальной модели для проверки H1 - максимум F1 среди конфигураций, не выходящих за полосу <<+150~\% L1-реконструкции от baseline>>. Указанное послабление полосы $L_1$ обусловлено тем, что критерий H1 формулируется не на $L_1$, а на LSE-C тестовой выборки; полоса $L_1$ выступает санитарным фильтром против заведомо испорченных конфигураций, а финальный вердикт по сохранности артикуляции выносится по LSE-C на отложенной тестовой выборке.

После отбора финальной модели проводится сравнение с базовой моделью (\texttt{wav2lip-baseline}, нулевой эмоциональный лосс) на тестовой выборке RAVDESS ($n=96$, по 24 сэмпла на эмоцию):
\begin{itemize}
\item категориальная метрика эмоциональной согласованности по внутреннему видеоэнкодеру (accuracy, макро-F1, per-emotion P/R/F1, матрица ошибок);
\item LSE-C / LSE-D, вычисляемые предобученным lipsync-экспертом SyncNet;
\item L1-реконструкция как контрольная величина.
\end{itemize}
Статистическая значимость: критерий Макнемара для F1 (с подсчётом дискордантных пар $n_{01}$, $n_{10}$, $\chi^2$, $p$); парный $t$-критерий и критерий Уилкоксона для L1 и LSE-C. Решение по H1 принимается при одновременном выполнении двух критериев: $\Delta\text{F1} \geq +10$~\% при $p < 0{,}05$ (Макнемар) и $|\Delta\text{LSE-C}| \leq 2$~\% либо $p \geq 0{,}05$ (парный $t$-критерий).

\subsection{Внешняя валидация и контроль цикличности оценки (этап 4, 5)}

Поскольку и функция потерь, и первичная метрика на этапе 5 опираются на один и тот же видеоэнкодер $\phi_v$, оценка только по $\phi_v$ создаёт риск кругового рассуждения. Для устранения риска проведён скрининг внешних классификаторов выражений лица в [external\_screening.ipynb](external\_screening.ipynb): из шести кандидатов один (\texttt{prithivMLmods/Facial-Emotion-Detection-SigLIP2}) исключён из-за отсутствия класса \textit{disgust} в его выходном пространстве; пять оставшихся прогнаны на тестовой выборке RAVDESS в режиме классификации центральных кадров реальных видео, и для основной внешней оценки отобраны два классификатора с максимальной макро-F1 на четырёх целевых эмоциях - \texttt{motheecreator/vit-Facial-Expression-Recognition} и \texttt{Rajaram1996/FacialEmoRecog}; \texttt{trpakov/vit-face-expression} удерживается как контрольный классификатор для отдельных подвыборок.

На этапе 5 финальные чекпойнты Wav2Lip применяются к тестовому подмножеству, генерируется выход на каждый сэмпл, центральный кадр из нижней половины каждого сгенерированного видео классифицируется каждым внешним классификатором, и подсчитывается макро-F1, accuracy, per-emotion F1 и матрица ошибок. Для каждой пары <<конфигурация, внешний классификатор>> рассчитывается $\Delta$F1 относительно базовой модели и проводится критерий Макнемара. В качестве диагностической метрики фиксируется \textit{cross-external agreement} - доля конфигураций, для которых направление $\Delta$F1 совпадает между двумя различающимися по архитектуре и обучающему корпусу классификаторами.

В качестве вспомогательного аудио-канала проверки фиксируется регрессия arousal/dominance/valence по модели \texttt{audeering/wav2vec2-large-robust-12-ft-emotion-msp-dim} (MSP-Podcast), что позволяет соотнести предсказанные эмоциональные категории с непрерывными размерными признаками речи и тем самым отделить случайные совпадения меток от содержательно согласованного эмоционального тона.

\subsection{Перенос на SadTalker и проверка H2 (этап 5)}

Опорная архитектура SadTalker [Zhang et al., 2023] отличается от Wav2Lip тем, что генерирует параметры 3DMM-геометрии лица, а не пиксели нижней половины кадра. Для применения той же функции потерь \eqref{eq:totalloss} в её декомпозиции <<ExpNet $\to$ FaceRender>> промежуточный сигнал отбирается на этапе синтеза кадров: после рендеринга применяется тот же видеоэнкодер $\phi_v$ к сгенерированному кадру, и далее используются $\mathcal{L}_{\text{CE}}$ и $\mathcal{L}_{\text{KL}}$. Файнтюнинг проводится на тех же 544 обучающих сэмплах RAVDESS, что и Wav2Lip; сохраняется тот же опорный аудиоэнкодер. Для скрининга масштаба $s$ в [07\_finetune\_sadtalker.ipynb](07\_finetune\_sadtalker.ipynb) проверены конфигурации $s=0{,}25$ и $s=1{,}0$; в [07a\_sadtalker\_loss\_ablation.ipynb](07a\_sadtalker\_loss\_ablation.ipynb) дополнительно проверена конфигурация \textit{abl-ce-cos}, заявленная как лучшая по средней внешней $\Delta$F1 в эксперименте Wav2Lip (см. результаты).

Внешняя оценка SadTalker отличается от Wav2Lip следующим: вместо классификации центральных кадров библиотечных тестовых видео внешний классификатор применяется к рендеренным видео (один кадр из каждых четырёх) с агрегированием по голосованию большинства; объём финальной выборки - 24 видео на конфигурацию (по 6 на эмоцию), что обусловлено вычислительной стоимостью рендеринга. Указанное ограничение явно отражено в разделе результатов и в обсуждении ограничений; малый размер выборки исключает чувствительное обнаружение слабых эффектов.

Решение по H2: $\Delta\text{F1} \geq +8$~\% по обоим внешним классификаторам с положительным согласованием направления и $p < 0{,}05$ по критерию Макнемара хотя бы для одного из них.

\subsection{Воспроизводимость}

Все ноутбуки публично доступны в репозитории (\url{https://github.com/Katrin-Pochtar/The-Uncanny-Valley}); фиксированный seed=42 для всех численных операций (Python \texttt{random}, NumPy, PyTorch, CUDA); версии библиотек закреплены в \texttt{requirements} и в выводе \texttt{pip list} первой ячейки каждого ноутбука; гиперпараметры и метрики записаны в проекты W\&B (\texttt{uncanny-valley-encoders-4emo}, \texttt{uncanny-valley-wav2lip-ablation}, \texttt{uncanny-valley-wav2lip}, \texttt{uncanny-valley-sadtalker}). Чекпойнты опорных энкодеров, абляции и финальных моделей сохранены в локальных каталогах, структура которых описана в README репозитория.

\section{Результаты}

\subsection{Подбор опорных энкодеров эмоций}

В рамках этапа~2 проведено 22 запуска (11 аудиоконфигураций и 11 видеоконфигураций); сводная таблица лучших валидационных макро-F1 представлена в табл.~\ref{tab:enc_sweep}.

\begin{table}[H]
\caption{Лучшие конфигурации энкодеров эмоций (валидационная макро-F1, RAVDESS, 4 эмоции)}
\label{tab:enc_sweep}
\centering
\small
\begin{tabular}{|p{5.5cm}|c|c|c|}
\hline
\textbf{Конфигурация} & \textbf{Модальность} & \textbf{Лучшая Val~F1} & \textbf{Финал} \\
\hline
\texttt{wav2vec2-base-superb-er}, lr $3\!\cdot\!10^{-5}$ & аудио & 0,826 & да \\
\hline
\texttt{hubert-large-ll60k}, lr $2\!\cdot\!10^{-5}$ & аудио & 0,804 & - \\
\hline
\texttt{wav2vec2-large}, lr $2\!\cdot\!10^{-5}$ & аудио & 0,801 & - \\
\hline
\texttt{timesformer-base}, 8 кадров, lr $3\!\cdot\!10^{-5}$ & видео & 0,790 & да \\
\hline
\texttt{timesformer-base}, 16 кадров, no-freeze, lr $3\!\cdot\!10^{-5}$ & видео & 0,781 & - \\
\hline
\texttt{timesformer-base}, 16 кадров, lr $5\!\cdot\!10^{-5}$ & видео & 0,780 & - \\
\hline
\end{tabular}
\end{table}

На независимой тестовой выборке отобранные энкодеры демонстрируют макро-F1 0,739 (аудио) и 0,876 (видео) ([04b\_encoder\_ceiling.ipynb](04b\_encoder\_ceiling.ipynb)). Поэмоциональная разбивка приведена в табл.~\ref{tab:enc_ceiling}; видеоэнкодер уверенно классифицирует \textit{happy}, \textit{angry} и \textit{disgust} (F1 $\geq$ 0,90) и хуже работает на \textit{sad} (F1 0,72), что согласуется с известной из литературы низкой различимостью <<печали>> по статичной мимике.

\begin{table}[H]
\caption{<<Потолок>> разделимости энкодеров: per-emotion F1 на тестовой выборке RAVDESS ($n=96$)}
\label{tab:enc_ceiling}
\centering
\small
\begin{tabular}{|p{4cm}|c|c|c|c|c|}
\hline
\textbf{Энкодер} & \textbf{Macro F1} & \textbf{F1 happy} & \textbf{F1 sad} & \textbf{F1 angry} & \textbf{F1 disgust} \\
\hline
Аудио (\texttt{wav2vec2-base-superb-er}) & 0,739 & 0,667 & 0,760 & 0,721 & 0,810 \\
\hline
Видео (\texttt{timesformer-base}, 8 fr.) & 0,876 & 0,906 & 0,718 & 0,923 & 0,958 \\
\hline
\end{tabular}
\end{table}

\subsection{Скрининг внешних классификаторов}

В рамках этапа~4 шесть кандидатов прогнаны на тестовой выборке RAVDESS ([external\_screening.ipynb](external\_screening.ipynb)); один (\texttt{prithivMLmods/Facial-Emotion-Detection-SigLIP2}) исключён из-за отсутствия класса \textit{disgust}. Результаты приведены в табл.~\ref{tab:ext_screen}.

\begin{table}[H]
\caption{Скрининг внешних классификаторов выражений лица на тестовой выборке RAVDESS, $n=96$}
\label{tab:ext_screen}
\centering
\small
\begin{tabular}{|p{6cm}|c|c|c|c|c|}
\hline
\textbf{Классификатор} & \textbf{Macro F1} & \textbf{F1 happy} & \textbf{F1 sad} & \textbf{F1 angry} & \textbf{F1 disgust} \\
\hline
\texttt{motheecreator/vit-FER} & \textbf{0,476} & 0,738 & 0,133 & 0,410 & 0,621 \\
\hline
\texttt{trpakov/vit-face-expression} & 0,464 & 0,814 & 0,364 & 0,677 & 0,000 \\
\hline
\texttt{Rajaram1996/FacialEmoRecog} & \textbf{0,462} & 0,744 & 0,250 & 0,263 & 0,592 \\
\hline
\texttt{RickyIG/emotion\_face} & 0,432 & 0,718 & 0,486 & 0,000 & 0,522 \\
\hline
\texttt{dima806/facial\_emotions} & 0,280 & 0,608 & 0,065 & 0,448 & 0,000 \\
\hline
\end{tabular}
\end{table}

Промежуточная макро-F1 внешних классификаторов на чистом RAVDESS (0,28-0,48) существенно ниже <<потолка>> внутреннего видеоэнкодера (0,876), что фиксирует эффект доменного сдвига между обучающими корпусами внешних классификаторов (преимущественно AffectNet/FER+) и студийной актёрской записью RAVDESS. Указанный сдвиг определяет верхнюю границу любого эффекта, регистрируемого внешними классификаторами на сгенерированном видео: даже идеальный генератор не может превзойти внешнюю F1 на реальных кадрах. Для финальной внешней оценки отобраны \texttt{motheecreator/vit-FER} и \texttt{Rajaram1996/FacialEmoRecog}.

\subsection{Абляция функции потерь (этап 3)}

Результаты пятиконфигурационной абляции на валидационной выборке Wav2Lip ($n=96$) представлены в табл.~\ref{tab:ablation}. Победившая по F1 конфигурация - \textbf{abl-ce-kl} (комбинация перекрёстной энтропии и KL-дистилляции).

\begin{table}[H]
\caption{Абляция функции потерь на Wav2Lip, валидационная выборка RAVDESS}
\label{tab:ablation}
\centering
\small
\begin{tabular}{|p{2.5cm}|c|c|c|c|c|c|}
\hline
\textbf{Конфигурация} & $w_{\text{ce}}$ & $w_{\cos}$ & $w_{\text{KL}}$ & \textbf{Best L1} & \textbf{Best F1} & $\Delta$F1 \\
\hline
abl-baseline & 0,00 & 0,0 & 0,00 & 0,0388 & 0,619 & 0,000 \\
\hline
abl-cos-only & 0,00 & 0,2 & 0,00 & 0,0420 & 0,574 & $-0,045$ \\
\hline
abl-ce-only & 0,05 & 0,0 & 0,00 & 0,0619 & 0,696 & $+0,077$ \\
\hline
abl-ce-cos & 0,05 & 0,2 & 0,00 & 0,0654 & 0,655 & $+0,036$ \\
\hline
\textbf{abl-ce-kl} & \textbf{0,05} & \textbf{0,0} & \textbf{0,10} & \textbf{0,0831} & \textbf{0,737} & \textbf{$+0,118$} \\
\hline
\end{tabular}
\end{table}

\textit{Содержательный вывод абляции:} наивная косинусная формулировка без явных негативных пар (abl-cos-only) приводит к \textit{отрицательной} $\Delta$F1 относительно baseline, что эмпирически воспроизводит описанный в литературе риск тривиального коллапса проекций [Chen, He, 2021]. Включение явной супервизии в виде перекрёстной энтропии (abl-ce-only) даёт умеренный прирост; добавление к ней KL-дистилляции от аудиоэнкодера к видеоэнкодеру (abl-ce-kl) увеличивает прирост ещё в полтора раза, при этом критерий Макнемара на $n=96$ показывает $p$-значения, не достигающие порога $p<0{,}05$ (для лучшей конфигурации abl-ce-kl: $n_{10}=22$, $n_{01}=14$, $p = 0{,}243$), что обусловлено малым объёмом валидационной выборки. Деградация L1-реконструкции значима ($p < 10^{-9}$ во всех эмоциональных конфигурациях): эмоциональный сигнал смещает декодер в сторону более эмоционально выраженной мимики ценой более грубой попиксельной реконструкции, что и составляет содержание исследуемого компромисса. Финальный вердикт по сохранности артикуляции выносится на тестовой выборке по LSE-C, не по L1.

Внешняя оценка той же абляции на тестовой выборке (этап 4, [05\_external\_evaluation.ipynb](05\_external\_evaluation.ipynb)) воспроизводит ту же иерархию: abl-cos-only имеет наименьший $\Delta$F1 среди эмоциональных конфигураций, abl-ce-cos и abl-ce-kl - максимальный; направление эффекта согласовано между двумя независимыми внешними классификаторами для всех четырёх эмоциональных конфигураций (cross-external agreement = True).

\subsection{H1: проверка на тестовой выборке Wav2Lip}

Победившая в абляции комбинация (CE+KL, $\tau=2{,}0$) применена к Wav2Lip с переменным масштабом $s$ ([06\_finetune\_wav2lip.ipynb](06\_finetune\_wav2lip.ipynb)). Сводная таблица валидационных результатов:

\begin{table}[H]
\caption{Поиск масштаба $s$ для CE+KL-функции потерь, валидационная выборка}
\label{tab:scale_sweep}
\centering
\small
\begin{tabular}{|p{2.7cm}|c|c|c|c|c|}
\hline
\textbf{Конфигурация} & $s$ & $w_{\text{ce}}$ & $w_{\text{KL}}$ & \textbf{Best L1} & \textbf{Best F1} \\
\hline
wav2lip-baseline & 0,00 & 0,0000 & 0,000 & 0,0382 & 0,624 \\
\hline
\textbf{wav2lip-cekl-01} & \textbf{0,10} & \textbf{0,0050} & \textbf{0,010} & \textbf{0,0483} & \textbf{0,742} \\
\hline
wav2lip-cekl-04 & 0,40 & 0,0200 & 0,040 & 0,0615 & 0,739 \\
\hline
wav2lip-cekl-05 & 0,50 & 0,0250 & 0,050 & 0,0925 & 0,736 \\
\hline
wav2lip-cekl-025 & 0,25 & 0,0125 & 0,025 & 0,0894 & 0,691 \\
\hline
\end{tabular}
\end{table}

Финальный чекпойнт - \textbf{wav2lip-cekl-01} (максимум валидационной F1 при наименьшей деградации L1; находится в полосе допустимой реконструкции). Сравнение с базовой моделью на тестовой выборке ($n=96$):

\begin{table}[H]
\caption{H1: \textit{wav2lip-baseline} vs \textit{wav2lip-cekl-01}, тестовая выборка RAVDESS}
\label{tab:h1_test}
\centering
\small
\begin{tabular}{|p{4cm}|c|c|c|c|}
\hline
\textbf{Метрика} & \textbf{Baseline} & \textbf{CE+KL} & $\Delta$ & $p$-value \\
\hline
Accuracy (внутр. видеоэнкодер) & 0,635 & 0,719 & $+0,083$ & 0,170 (Mc) \\
\hline
Macro F1 (внутр. видеоэнкодер) & 0,611 & 0,718 & $+0,107$ & 0,170 (Mc) \\
\hline
LSE-C $\uparrow$ & 0,242 & 0,227 & $-0,015$ ($-6{,}1$~\%) & 0,611 ($t$) \\
\hline
LSE-D $\downarrow$ & 0,758 & 0,773 & $+0,015$ & 0,611 ($t$) \\
\hline
L1-реконструкция $\downarrow$ & 0,047 & 0,072 & $+0,025$ & $4{,}4\!\cdot\!10^{-13}$ ($W$) \\
\hline
\end{tabular}
\end{table}

Поэмоциональная разбивка F1 (табл.~\ref{tab:h1_per_emo}) показывает, что наибольший прирост приходится на классы, наиболее <<пересекаемые>> с базовой моделью - \textit{sad} ($+0{,}205$) и \textit{angry} ($+0{,}120$); эффект на \textit{happy} мал ($+0{,}031$) ввиду насыщения базовой модели на этом классе (recall = 1,000).

\begin{table}[H]
\caption{H1: per-emotion F1 на тестовой выборке (Wav2Lip)}
\label{tab:h1_per_emo}
\centering
\small
\begin{tabular}{|c|c|c|c|}
\hline
\textbf{Эмоция} & \textbf{Baseline F1} & \textbf{CE+KL F1} & $\Delta$F1 \\
\hline
happy & 0,762 & 0,792 & $+0,031$ \\
\hline
sad & 0,378 & 0,583 & $+0,205$ \\
\hline
angry & 0,653 & 0,773 & $+0,120$ \\
\hline
disgust & 0,651 & 0,723 & $+0,072$ \\
\hline
\end{tabular}
\end{table}

\textit{Вердикт по H1.} Магнитуда эффекта по F1 удовлетворяет порогу гипотезы ($\Delta\text{F1} = +10{,}7$~\% $\geq +10$~\%); статистическая значимость по критерию Макнемара на $n=96$ не достигнута ($p = 0{,}170$). Изменение LSE-C по абсолютной величине превосходит порог в 2~\% ($\Delta = -6{,}1$~\%), однако не является статистически значимым ($p = 0{,}611$). Согласно сформулированному в методологии правилу решения, $|\Delta\text{LSE-C}|>2$~\% при $p \geq 0{,}05$ удовлетворяет ослабленному критерию (изменение неотличимо от нуля); таким образом, ограничение на сохранность артикуляции \textbf{не нарушено}, тогда как улучшение F1 \textbf{не получило статистического подтверждения} на текущем размере выборки. \textbf{Гипотеза H1 в её строгой формулировке не отвергнута и не подтверждена}; направление эффекта систематично (одинаково на двух внешних классификаторах, см. ниже), но статистическая мощность $n=96$ недостаточна для обнаружения эффекта искомой величины.

\subsection{H1: внешняя валидация}

Внешняя оценка той же пары моделей на двух независимых классификаторах (табл.~\ref{tab:h1_external}) показывает положительный $\Delta$F1 на обоих и направленную согласованность между внешними классификаторами и внутренней метрикой.

\begin{table}[H]
\caption{H1: внешняя валидация на тестовой выборке (Wav2Lip)}
\label{tab:h1_external}
\centering
\small
\begin{tabular}{|p{3.5cm}|c|c|c|c|}
\hline
\textbf{Внешний классификатор} & \textbf{Baseline F1} & \textbf{CE+KL F1} & $\Delta$F1 & McNemar $p$ \\
\hline
\texttt{motheecreator/vit-FER} & 0,303 & 0,343 & $+0,040$ & 0,286 \\
\hline
\texttt{Rajaram1996/FacialEmoRecog} & 0,338 & 0,344 & $+0,007$ & 1,000 \\
\hline
\end{tabular}
\end{table}

Cross-external agreement по направлению эффекта - True для финального чекпойнта \textit{cekl-01}. На альтернативных конфигурациях из той же CE+KL-сем\'ьи достигаются б\'ольшие положительные $\Delta$F1 на обоих внешних классификаторах (\textit{cekl-04}: $+0{,}059$ и $+0{,}039$; \textit{cekl-05}: $+0{,}033$ и $+0{,}089$ соответственно), но при б\'ольшей деградации $L_1$-реконструкции, что и обусловило выбор \textit{cekl-01} как финальной модели по правилу селекции <<максимум F1 при минимуме L1>>. Статистическая значимость на $n=96$ не достигнута ни для одной пары; направление эффекта при этом устойчиво между двумя внешними классификаторами, что снижает риск артефакта одного источника оценки.

\subsection{H2: проверка на SadTalker}

Опорная архитектура SadTalker файнтюнирована с применением CE+KL-комбинации при $s=0{,}25$ и $s=1{,}0$ ([07\_finetune\_sadtalker.ipynb](07\_finetune\_sadtalker.ipynb)) и дополнительно с конфигурацией \textit{abl-ce-cos} ([07a\_sadtalker\_loss\_ablation.ipynb](07a\_sadtalker\_loss\_ablation.ipynb)). Внешняя оценка по рендеренным видео ($n=24$ на конфигурацию, по 6 на эмоцию) представлена в табл.~\ref{tab:h2_test}.

\begin{table}[H]
\caption{H2: внешняя оценка SadTalker на рендеренных тестовых видео ($n=24$)}
\label{tab:h2_test}
\centering
\small
\begin{tabular}{|p{4cm}|c|c|c|c|}
\hline
\textbf{Конфигурация} & \multicolumn{2}{c|}{\texttt{trpakov} / \texttt{motheecreator}} & \multicolumn{2}{c|}{\texttt{Rajaram1996} / \texttt{dima806}} \\
 & \textbf{F1} & $\Delta$F1 & \textbf{F1} & $\Delta$F1 \\
\hline
sadtalker-baseline (07) & 0,383 (\textit{trpakov}) & 0,000 & 0,215 (\textit{dima806}) & 0,000 \\
\hline
sadtalker-cekl-025 & 0,369 & $-0,015$ & 0,215 & 0,000 \\
\hline
sadtalker-cekl-10 & 0,371 & $-0,012$ & 0,215 & 0,000 \\
\hline
abl-baseline (07a) & 0,409 (\textit{motheecreator}) & 0,000 & 0,373 (\textit{Rajaram1996}) & 0,000 \\
\hline
abl-ce-cos (07a) & 0,467 & $+0,058$ & 0,181 & $-0,193$ \\
\hline
\end{tabular}
\end{table}

\textit{Вердикт по H2.} На рендеренной выборке SadTalker не зафиксирован устойчивый положительный эффект, согласованный между двумя независимыми внешними классификаторами: \textit{cekl} на \textit{trpakov} даёт малое отрицательное смещение (статистически незначимое, McNemar $p=1{,}0$ при $n_{01}=2$, $n_{10}=3$); \textit{abl-ce-cos} на \textit{motheecreator} показывает положительный $\Delta$F1, однако на \textit{Rajaram1996} - крупное отрицательное смещение (cross-external agreement = False). Критерий Макнемара ни на одной паре не достигает $p<0{,}05$. \textbf{Гипотеза H2 в текущем эксперименте не подтверждена.} Возможные причины и ограничения обсуждаются в следующем подразделе.

\subsection{Ограничения проведённого эксперимента}

Перечисленные ниже ограничения определяют границы применимости полученных результатов и не должны игнорироваться при их интерпретации:
\begin{enumerate}
\item \textbf{Размер тестовых выборок.} Тестовая выборка Wav2Lip составляет 96 сэмплов, тестовая выборка SadTalker по рендеренным видео - 24 сэмпла на конфигурацию. Указанные размеры обусловлены актёрно-разделённым разбиением RAVDESS (24 актёра $\to$ 2-3 актёра на сплит) и вычислительной стоимостью полного рендеринга SadTalker. Малая выборка существенно ограничивает статистическую мощность критерия Макнемара: для обнаружения эффекта величиной $\Delta\text{F1}=+0{,}1$ при достижимом дискордансе требуется $n \gtrsim 200$, что объясняет наблюдаемые $p$-значения порядка 0,15-0,30 при содержательно интерпретируемой магнитуде эффекта.
\item \textbf{Студийный характер RAVDESS.} 24 актёра, две стандартные фразы, фиксированная поза головы, контролируемое освещение - параметры, не воспроизводимые в in-the-wild сценариях.
\item \textbf{Сужение эмоционального пространства до 4 классов.} Исключение neutral/calm/fearful/surprised мотивировано в теоретической части; обобщение на полное восьмиклассовое пространство требует отдельного эксперимента.
\item \textbf{Доменный сдвиг внешних классификаторов.} Макро-F1 внешних классификаторов на чистых кадрах RAVDESS составляет 0,28-0,48 против 0,876 у внутреннего видеоэнкодера. Это создаёт <<потолок>> внешней оценки и снижает чувствительность к слабым эффектам.
\item \textbf{Архитектурный перенос на SadTalker.} Текущая интеграция функции потерь применяется к рендерному выходу SadTalker; альтернативная интеграция на уровне промежуточных 3DMM-параметров (ExpNet) в рамках настоящего исследования не проверена и составляет одно из направлений дальнейших работ.
\item \textbf{Отсутствие перцепционной оценки.} Объектные метрики (внутренняя F1, внешняя F1, LSE-C) не заменяют слепое A/B тестирование с участием респондентов, выходящее за пределы настоящего эксперимента.
\end{enumerate}

\section{Выводы}

\begin{enumerate}
\item \textbf{Гипотеза H1} (Wav2Lip): магнитуда наблюдаемого эффекта удовлетворяет заявленному порогу ($\Delta\text{F1} = +10{,}7$~\% при пороге $\geq +10$~\%) при сохранении не значимого по парному $t$-критерию изменения LSE-C; статистическая значимость прироста F1 по критерию Макнемара на $n=96$ не достигнута ($p = 0{,}170$). Гипотеза в строгой формулировке \textbf{не отвергнута и не подтверждена}; полученное направление эффекта систематично, согласовано между внутренней метрикой и двумя независимыми внешними классификаторами.
\item \textbf{Гипотеза H2} (SadTalker): на рендеренной выборке $n=24$ устойчивого положительного эффекта, согласованного между двумя внешними классификаторами, не зафиксировано; гипотеза \textbf{не подтверждена} в текущем эксперименте.
\item Абляционный анализ на пяти конфигурациях функции потерь показал, что наивная косинусная формулировка без явных негативных пар приводит к отрицательной $\Delta$F1 (коллапс проекций), тогда как комбинация перекрёстной энтропии и KL-дистилляции (CE+KL) даёт максимальный прирост; обоснованность выбора компонент эмпирически подтверждена.
\item Установлен <<потолок>> разделимости опорных энкодеров на тестовой выборке (макро-F1 0,876 для видеоэнкодера; 0,739 для аудиоэнкодера), задающий теоретическую верхнюю границу любого эффекта, оцениваемого через указанные энкодеры.
\item По итогам скрининга пяти внешних классификаторов выражений лица отобраны два (\texttt{motheecreator/vit-FER}, \texttt{Rajaram1996/FacialEmoRecog}), не пересекающиеся ни с обучающей сигнальной цепью, ни с метриками отбора. Их применение в качестве независимого валидатора устраняет круговое рассуждение, присущее значительной части работ по эмоциональному синтезу говорящих лиц.
\item Заявленный протокол реализации plug-in функции потерь к опорным архитектурам без модификации backbone и без явных меток на инференсе технически работоспособен и в контексте Wav2Lip даёт направленный положительный эффект; строгое статистическое подтверждение требует расширения тестовой выборки и/или применения корпусов с большим числом субъектов (CREMA-D, MEAD).
\end{enumerate}

\section{Перспективы и рекомендации}

Перечисленные ниже направления формулируются как ответы на конкретные ограничения, изложенные в подразделе ограничений; каждое направление обосновано полученными в работе данными.

\begin{enumerate}
\item \textbf{Расширение тестовой выборки.} Прямой ответ на ограничение статистической мощности: проведение того же эксперимента на корпусе CREMA-D ($n_{\text{актёров}}=91$, $n_{\text{записей}}=7442$) и MEAD ($n_{\text{актёров}}=60$ с тремя уровнями интенсивности) обеспечит размер тестового сплита $\gtrsim 500$, при котором критерий Макнемара получит мощность для обнаружения эффекта порядка $\Delta\text{F1}=+0{,}1$ при $p<0{,}05$.
\item \textbf{Расширение эмоционального пространства.} Переход от четырёхклассовой постановки к полному восьмиклассовому пространству RAVDESS либо к восьмиклассовому пространству MEAD с тремя уровнями интенсивности позволит проверить переносимость найденной CE+KL-комбинации на классы с более высокой межклассовой конфузией (\textit{fearful}/\textit{surprised}).
\item \textbf{Альтернативная интеграция функции потерь в SadTalker.} Текущая интеграция действует на рендерном выходе; перспективой является применение CE+KL-сигнала непосредственно к промежуточным 3DMM-параметрам выражений (выходы ExpNet) до этапа рендеринга, что снимет ограничение, обусловленное стохастичностью FaceRender и снижающее чувствительность внешней оценки.
\item \textbf{Перцепционная валидация.} Слепое A/B тестирование с группой респондентов на парах <<baseline / CE+KL>> позволит проверить, переносится ли регистрируемое автоматическими классификаторами улучшение в субъективное восприятие реализма аватара.
\item \textbf{Внешняя оценка в условиях in-the-wild.} Применение того же протокола к видео без студийного контроля (например, к подмножеству VoxCeleb с автоматической эмоциональной разметкой) позволит оценить устойчивость эффекта к доменному сдвигу.
\item \textbf{Внешняя апробация.} Основные результаты исследования представлены на 68-й Всероссийской научной конференции МФТИ (секция когнитивных технологий, 4 апреля 2026 г.), что выступает первым этапом внешней верификации; следующий шаг - подача расширенной статьи в специализированный рецензируемый журнал по теме мультимодальной генерации с включением расширенного эксперимента на CREMA-D.
\item \textbf{Прикладные рекомендации.} Для разработчиков, эксплуатирующих опорные архитектуры Wav2Lip-семейства для прикладных аватаров, продемонстрированный CE+KL-сигнал может быть применён без модификации архитектуры в режиме файнтюнинга на доменно-релевантном эмоциональном корпусе при ожидаемой деградации L1-реконструкции порядка $+50$~\% и ожидаемом сохранении LSE-C в пределах $\pm 6$~\%, что приемлемо для большинства целевых применений (цифровые ассистенты, образовательные платформы) и требует осторожной валидации в задачах, чувствительных к точности артикуляции (медиа-производство, дубляж).
\end{enumerate}

\section{Заключение}

В диссертации сформулирована и экспериментально проверена постановка задачи повышения эмоциональной согласованности аудио-управляемых моделей генерации говорящих лиц через plug-in кросс-модальную функцию потерь, не требующую модификации опорной архитектуры и явной разметки эмоций на этапе инференса. Теоретический анализ систематизирует существующие подходы к эмоциональному синтезу и выявляет лакуну, состоящую в отсутствии работ, в которых эмоциональная согласованность достигается надстройкой над уже развёрнутыми моделями с независимой внешней верификацией. Методическая часть формализует функцию потерь как сумму реконструкционного и кросс-модального членов, состав которых эмпирически устанавливается в пятиконфигурационной абляции, а гиперпараметры опорных энкодеров отбираются перебором из 22 конфигураций.

Эмпирическая часть фиксирует следующие результаты. Победившая в абляции комбинация \textit{CE+KL} применима к Wav2Lip и даёт на тестовой выборке прирост макро-F1 эмоциональной согласованности с 0,611 до 0,718 ($+10{,}7$~\%) при изменении LSE-C на $-6{,}1$~\% (статистически незначимом). Магнитуда эффекта удовлетворяет порогу гипотезы H1, однако статистическая значимость по критерию Макнемара на $n=96$ не достигнута, что обусловлено ограничением размера актёрно-разделённой тестовой выборки RAVDESS. Направление эффекта подтверждается двумя независимыми внешними классификаторами, отобранными по итогам скрининга, что устраняет риск кругового рассуждения. На SadTalker та же функция потерь в текущей интеграции не показала устойчивого согласованного эффекта; гипотеза H2 не подтверждена. Полные ограничения работы - размер выборок, студийный характер корпуса, четырёхклассовая постановка, доменный сдвиг внешних классификаторов и текущая интеграция в SadTalker - явно перечислены в соответствующем подразделе.

Дополнительные задачи, выходящие за рамки исследовательской части, включают: подготовку открытого репозитория с воспроизводимым пайплайном (\url{https://github.com/Katrin-Pochtar/The-Uncanny-Valley}); публикацию W\&B-проектов с полными журналами экспериментов; проведение бутстрап-оценки доверительных интервалов $\Delta$F1 ([bootstrap\_analysis.ipynb](bootstrap\_analysis.ipynb)). Основные результаты приняты к публикации в сборнике тезисов 68-й Всероссийской научной конференции МФТИ (\url{https://conf.mipt.ru/conference/23659}); устное представление состоялось 4 апреля 2026 г. Объекты ИС в рамках работы не регистрировались. Прикладное внедрение в эксплуатацию на момент сдачи работы не оформлено; рекомендации к внедрению приведены в разделе перспектив.

% \begin{landscape}
% % текст в альбомной ориентации
% % (таблица, рисунок, схема и т. п.)
% \end{landscape}

\addcontentsline{toc}{section}{Список литературы}
\nocite{*}  % ✅ Показать ВСЕ записи из .bib, даже без цитирований
\bibliographystyle{utf8gost705u}
\bibliography{my_bibliography}
%\printbibliography[title={Список литературы}] % выводит список литературы
%\nocite{*}  % ✅ Показать ВСЕ записи из .bib, даже без цитирований

%% Если не хочется заморачиваться с отдельным библиотечным файлом:
% \begin{thebibliography}{9} 
% \bibitem{lamport94} Leslie Lamport, \emph{\LaTeX: a document preparation system}. Addison Wesley, Massachusetts, 2nd edition, 1994. 
% \bibitem{lamport941} Leslie Lamport, \emph{\LaTeX: a document preparation system}. Addison Wesley, Massachusetts, 2nd edition, 1994. 
% \end{thebibliography}

\newpage

\appendix %Тут идут приложения
\renewcommand{\thesection}{Приложение \Asbuk{section}}
\section{\centering}\label{PrilA}
\begin{center} 
\hspace*{-7mm}{\large\textbf{Артефакты}}
\end{center}
\setcounter{equation}{0} % Сбрасываем счетчик формул
\renewcommand{\theequation}{\Asbuk{section}.\arabic{equation}} % Устанавливаем префикс "А."
\setcounter{table}{0} % Сбрасываем счетчик формул
\renewcommand{\thetable}{\Asbuk{section}.\arabic{table}} % Устанавливаем префикс "А." для таблицы
\setcounter{figure}{0} % Сбрасываем счетчик формул
\renewcommand{\thefigure}{\Asbuk{section}.\arabic{figure}} % Устанавливаем префикс "А." для рисунка

Все артефакты исследования открыто доступны и связаны с конкретными шагами методики и конкретными результатами; отсылка <<Артефакт $\to$ Шаг методики $\to$ Результат>> приведена ниже.

\begin{table}[H]
\caption{Реестр артефактов исследования}
\label{tab:artifacts}
\centering
\small
\begin{tabular}{|p{4.7cm}|p{3.5cm}|p{6.5cm}|}
\hline
\textbf{Артефакт} & \textbf{Шаг методики} & \textbf{Конкретный результат} \\
\hline
[01\_data\_preprocessing.ipynb](01\_data\_preprocessing.ipynb) & Подготовка данных & Сбалансированные сплиты 544/96/96 без пересечения по актёрам \\
\hline
[02\_train\_encoders\_4emotions.ipynb](02\_train\_encoders\_4emotions.ipynb) & Файнтюнинг энкодеров & Аудиоэнкодер val~F1=0,826; видеоэнкодер val~F1=0,790 \\
\hline
[03\_emotion\_utils.ipynb](03\_emotion\_utils.ipynb) & Утилиты загрузки/инференса энкодеров & Унификация интерфейса для последующих этапов \\
\hline
[04a\_wav2lip\_loss\_ablation.ipynb](04a\_wav2lip\_loss\_ablation.ipynb) & Абляция функции потерь & Победившая конфигурация CE+KL ($\Delta$F1=$+0{,}118$) \\
\hline
[04b\_encoder\_ceiling.ipynb](04b\_encoder\_ceiling.ipynb) & <<Потолок>> энкодеров & Test macro-F1: аудио 0,739; видео 0,876 \\
\hline
[external\_screening.ipynb](external\_screening.ipynb) & Скрининг внешних классификаторов & Отбор \texttt{motheecreator/vit-FER}, \texttt{Rajaram1996/FacialEmoRecog} \\
\hline
[05\_external\_evaluation.ipynb](05\_external\_evaluation.ipynb) & Внешняя валидация Wav2Lip & Cross-external agreement = True для всех CE+KL-конфигураций \\
\hline
[06\_finetune\_wav2lip.ipynb](06\_finetune\_wav2lip.ipynb) & Файнтюнинг Wav2Lip + проверка H1 & Финальная модель \textit{wav2lip-cekl-01}: F1 0,611$\to$0,718 \\
\hline
[07\_finetune\_sadtalker.ipynb](07\_finetune\_sadtalker.ipynb) & Файнтюнинг SadTalker + проверка H2 & Внешний $\Delta$F1$\approx 0$ на $n=24$; H2 не подтверждена \\
\hline
[07a\_sadtalker\_loss\_ablation.ipynb](07a\_sadtalker\_loss\_ablation.ipynb) & Дополнительная абляция SadTalker & Cross-external agreement = False для \textit{abl-ce-cos} \\
\hline
[bootstrap\_analysis.ipynb](bootstrap\_analysis.ipynb) & Бутстрап-анализ $\Delta$F1 & Доверительные интервалы $\Delta$F1 на парных бутстрап-выборках \\
\hline
W\&B \texttt{uncanny-valley-encoders-4emo} & Этап 2 & Полные журналы 22 запусков энкодеров \\
\hline
W\&B \texttt{uncanny-valley-wav2lip-ablation} & Этап 3 & Полные журналы 5 абляционных запусков \\
\hline
W\&B \texttt{uncanny-valley-wav2lip} & Этап 5 (Wav2Lip) & Полные журналы файнтюнинга и метрик \\
\hline
W\&B \texttt{uncanny-valley-sadtalker} & Этап 5 (SadTalker) & Полные журналы файнтюнинга и метрик \\
\hline
\url{https://github.com/Katrin-Pochtar/The-Uncanny-Valley} & Все этапы & Открытый репозиторий с воспроизводимым пайплайном \\
\hline
\end{tabular}
\end{table}


\newpage
\section{\centering}\label{PrilB}
\begin{center}
\hspace*{-7mm}{\large\textbf{Реализация кросс-модальной функции потерь}}
\end{center}
\setcounter{equation}{0}
\renewcommand{\theequation}{\Asbuk{section}.\arabic{equation}}
\setcounter{table}{0}
\renewcommand{\thetable}{\Asbuk{section}.\arabic{table}}
\setcounter{figure}{0}
\renewcommand{\thefigure}{\Asbuk{section}.\arabic{figure}}

Ниже приведена референсная реализация кросс-модальной функции потерь \eqref{eq:totalloss} в победившей по итогам абляции конфигурации CE+KL. Полная реализация с поддержкой всех пяти конфигураций абляции доступна в [04a\_wav2lip\_loss\_ablation.ipynb](04a\_wav2lip\_loss\_ablation.ipynb), производственная версия для финальных запусков - в [06\_finetune\_wav2lip.ipynb](06\_finetune\_wav2lip.ipynb).

\begin{lstlisting}[label=lst:loss, caption=Кросс-модальная функция потерь (CE+KL)]
import torch
import torch.nn.functional as F


def cross_modal_emotion_loss(
    pred_video, target_video,            # for L1 reconstruction
    audio_logits, video_logits,          # 4-dim emotion logits
    emotion_label,                       # int64, shape [B]
    s=0.10, w_ce=0.005, w_kl=0.010, T=2.0,
):
    # L1 over the lower half of the face (Wav2Lip-style)
    loss_recon = F.l1_loss(pred_video, target_video)

    # Cross-entropy on labels: explicit supervision
    loss_ce = F.cross_entropy(video_logits, emotion_label)

    # KL distillation: transfer the softened audio distribution
    # onto the video classifier head at temperature T
    student = F.log_softmax(video_logits / T, dim=-1)
    teacher = F.softmax(audio_logits.detach() / T, dim=-1)
    loss_kl = F.kl_div(student, teacher, reduction="batchmean") * (T ** 2)

    loss = loss_recon + s * (w_ce * loss_ce + w_kl * loss_kl)
    return loss, {"recon": loss_recon, "ce": loss_ce, "kl": loss_kl}
\end{lstlisting}

\textit{Соответствие \eqref{eq:totalloss}.} Параметры $\lambda_{\text{ce}}=0{,}05$, $\lambda_{\text{KL}}=0{,}1$, $\tau=2{,}0$ выбраны абляцией; масштаб $s=0{,}10$ выбран по результатам поиска веса в [06\_finetune\_wav2lip.ipynb](06\_finetune\_wav2lip.ipynb). Аудиологиты дистиллируются с \texttt{detach()}, поскольку аудиоэнкодер на этапе 5 заморожен и не должен принимать градиент. Косинусный член из \eqref{eq:totalloss} в финальной конфигурации обнуляется ($\lambda_{\cos}=0$) на основании отрицательного $\Delta$F1 наивной косинусной формулировки в абляции (см. табл.~\ref{tab:ablation}).
%\captionof{lstlisting}{Фрагмент кода функции subfinder}  % ✅ Русская подпись отдельно
% \begin{minted}[frame=lines, framesep=2mm, fontsize=\small, label=some-code,caption=Фрагмент кода функции subfinder]{python}
% def subfinder(words_list, answer_list): # code of function
%   matches = []
%   start_indices = []
%   end_indices = []
%   for idx, i in enumerate(range(len(words_list))):
%     if words_list[i] == answer_list[0] and words_list[i : i + len(answer_list)] == answer_list:
%       matches.append(answer_list)
%       start_indices.append(idx)
%       end_indices.append(idx + len(answer_list) - 1)
%     if matches:
%       return matches[0], start_indices[0], end_indices[0]
%     else:
%       return None, 0, 0
% \end{minted}

% Импорт из внешнего файла:
% \lstinputlisting[language=Python, frame=single]{script.py}



% \intro{Название приложения}\label{PrilA}


% Содержание приложения

% \newpage
% \intro{Название приложения}
% \label{PrilB}

% Содержание приложения

\end{document}

